% Generated by tools/make_body.py; edit generator, not this file.
% UNPROOFREAD OCR TRANSCRIPTION: page-image verification remains required.
\chapter{Series and Product Developments}\label{chap:5}

% Original PDF physical page 191
\sourcepage{191}
Very general theorems have their natural place in the theory of analytic functions, but it must also be kept in mind that the whole theory originated from a desire to be able to manipulate explicit analytic expressions. Such expressions take the form of infinite series, infinite products, and other limits. In this chapter we deal partly with the rules that govern such limits, partly with quite explicit representations of elementary transcendental functions and other specific functions.

\setcounter{section}{0}\section{Power Series Expansions}\label{sec:5.1}
In a preliminary way we have considered power series in Chap.~\hyperref[chap:2]{2}, mainly for the purpose of defining the exponential and trigonometric functions. Without use of integration we were not able to prove that every analytic function has a power series expansion. This question will now be resolved in the affirmative, essentially as an application of Cauchy's theorem.

The first subsection deals with more general properties of sequences of analytic functions.

\setcounter{section}{1}\setcounter{subsection}{0}\subsection{Weierstrass's Theorem}\index{Weierstrass's theorem|(}\index{Weierstrass's theorem}\label{sec:5.1.1}
The central theorem concerning the convergence of analytic functions asserts that the limit of a uniformly convergent sequence of analytic functions is an analytic function. The precise assumptions must be carefully stated, and they should not be too restrictive.


% Original PDF physical page 192
\sourcepage{192}
We are considering a sequence $\{f_{n}(z)\}$ where each $f_{n}(z)$ is defined and analytic in a region $\Omega_{n}$. The limit function $f(z)$ must also be considered in some region $\Omega$, and clearly, if $f(z)$ is to be defined in $\Omega$, each point of $\Omega$ must belong to all $\Omega_{n}$ for n greater than a certain $n_{0}$. In the general case $n_{0}$ will not be the same for all points of $\Omega$, and for this reason it would not make sense to require that the convergence be uniform in $\Omega$. In fact, in the most typical case the regions $\Omega_{n}$ form an increasing sequence, $\Omega_{1} \subset \Omega_{2} \subset \cdots \subset \Omega_{n} \subset \cdots$, and $\Omega$ is the union of the $\Omega_{n}$. In these circumstances no single function $f_{n}(z)$ is defined in all of $\Omega$; yet the limit $f(z)$ may exist at all points of $\Omega$, although the convergence cannot be uniform.

As a very simple example take $f_{n}(z) = z/(2z^{n} + 1)$ and let $\Omega_{n}$ be the disk $|z| < 2^{-1/n}$. It is practically evident that $\lim_{n \to \infty} f_{n}(z) = z$ in the disk $|z| < 1$ which we choose as our region $\Omega$. In order to study the uniformity of the convergence we form the difference
\begin{equation*}
f _ {n} (z) - z = - 2 z ^ {n + 1} / (2 z ^ {n} + 1).
\end{equation*}
For any given value of $z$ we can make $|z^n| < \varepsilon / 4$ by taking $n > \log (4 / \varepsilon) / \log (1 / |z|)$. If $\varepsilon < 1$ we have then $2|z|^{n+1} < \varepsilon / 2$ and $|1 + 2z^n| > \frac{1}{2}$ so that $|f_n(z) - z| < \varepsilon$. It follows that the convergence is uniform in any closed disk $|z| \leq r < 1$, or on any subset of such a closed disk.

With another formulation, in the preceding example the sequence $\{f_{n}(z)\}$ tends to the limit function $f(z)$ uniformly on every compact subset of the region $\Omega$. In fact, on a compact set $|z|$ has a maximum r < 1 and the set is thus contained in the closed disk $|z| \leq r$. This is the typical situation. We shall find that we can frequently prove uniform convergence on every compact subset of $\Omega$; on the other hand, this is the natural condition in the theorem that we are going to prove.

\begin{originalstatement}{Theorem}{1}\label{thm:5.1}
Suppose that $f_{n}(z)$ is analytic in the region $\Omega_{n}$, and that the sequence $\{f_{n}(z)\}$ converges to a limit function $f(z)$ in a region $\Omega$, uniformly on every compact subset of $\Omega$. Then $f(z)$ is analytic in $\Omega$. Moreover, $f_{n}'(z)$ converges uniformly to $f'(z)$ on every compact subset of $\Omega$.
\end{originalstatement}
The analyticity of $f(z)$ follows most easily by use of Morera's theorem (Chap.~\hyperref[chap:4]{4}, Sec.~\hyperref[sec:4.2.3]{2.3}). Let $|z - a| \leq r$ be a closed disk contained in $\Omega$; the assumption implies that this disk lies in $\Omega_n$ for all $n$ greater than a certain $n_0$.\hyperlink{fn:192:dagger}{\textsuperscript{\ensuremath{\dagger}}} If $\gamma$ is any closed curve contained in $|z - a| < r$, we have
\begin{equation*}
\int_ {\gamma} f _ {n} (z) d z = 0
\end{equation*}
\begingroup\renewcommand{\thefootnote}{}\footnotetext{\hypertarget{fn:192:dagger}{}\textsuperscript{\ensuremath{\dagger}} In fact, the regions $\Omega_{n}$ form an open covering of $|z - a| \leq r$. The disk is compact and hence has a finite subcovering. This means that it is contained in a fixed $\Omega_{n_{0}}$.}\endgroup

% Original PDF physical page 193
\sourcepage{193}
for $n > n_0$, by Cauchy's theorem. Because of the uniform convergence on $\gamma$ we obtain
\begin{equation*}
\int_ {\gamma} f (z) d z = \lim _ {n \rightarrow \infty} \int_ {\gamma} f _ {n} (z) d z = 0,
\end{equation*}
and by Morera's theorem it follows that $f(z)$ is analytic in $|z - a| < r$. Consequently $f(z)$ is analytic in the whole region $\Omega$.

An alternative and more explicit proof is based on the integral formula
\begin{equation*}
f _ {n} (z) = \frac {1}{2 \pi i} \int_ {C} \frac {f _ {n} (\zeta) d \zeta}{\zeta - z},
\end{equation*}
where $C$ is the circle $|\zeta - a| = r$ and $|z - a| < r$. Letting $n$ tend to $\infty$ we obtain by uniform convergence
\begin{equation*}
f (z) = \frac {1}{2 \pi i} \int_ {C} \frac {f (\zeta) d \zeta}{\zeta - z}.
\end{equation*}
and this formula shows that $f(z)$ is analytic in the disk. Starting from the formula
\begin{equation*}
f _ {n} ^ {\prime} (z) = \frac {1}{2 \pi i} \int_ {C} \frac {f _ {n} (\zeta) d \zeta}{(\zeta - z) ^ {2}}
\end{equation*}
the same reasoning yields
\begin{equation*}
\lim _ {n \rightarrow \infty} f _ {n} ^ {\prime} (z) = \frac {1}{2 \pi i} \int_ {C} \frac {f (\zeta) d \zeta}{(\zeta - z) ^ {2}} = f ^ {\prime} (z),
\end{equation*}
and simple estimates show that the convergence is uniform for $|z - a| \leq \rho < r$. Any compact subset of $\Omega$ can be covered by a finite number of such closed disks, and therefore the convergence is uniform on every compact subset. The theorem is proved, and by repeated applications it follows that $f_{n}^{(k)}(z)$ converges uniformly to $f^{(k)}(z)$ on every compact subset of $\Omega$.

Theorem~\hyperref[thm:5.1]{1} is due to Weierstrass, in an equivalent formulation. Its application to series whose terms are analytic functions is particularly important. The theorem can then be expressed as follows:

\begin{originalassertion}
If a series with analytic terms,
\begin{equation*}
f (z) = f _ {1} (z) + f _ {2} (z) + \dots + f _ {n} (z) + \dots ,
\end{equation*}
converges uniformly on every compact subset of a region $\Omega$, then the sum $f(z)$ is analytic in $\Omega$, and the series can be differentiated term by term.
\end{originalassertion}

The task of proving uniform convergence on a compact point set A can be facilitated by use of the maximum principle. In fact, with the notations of Theorem~\hyperref[thm:5.1]{1}, the difference $|f_{m}(z) - f_{n}(z)|$ attains its maxi%
% Original PDF physical page 194
\sourcepage{194}%
mum in $A$ on the boundary of $A$. For this reason uniform convergence on the boundary of $A$ implies uniform convergence on $A$. For instance, if the functions $f_n(z)$ are analytic in the disk $|z| < 1$, and if it can be shown that the sequence converges uniformly on each circle $|z| = r_m$ where $\lim_{m \to \infty} r_m = 1$, then Weierstrass's theorem\index{Weierstrass's theorem} applies and we can conclude that the limit function is analytic.

The following theorem is due to A. Hurwitz\index{Hurwitz A@Hurwitz, A.}:

\begin{originalstatement}{Theorem}{2}\label{thm:5.2}
If the functions $f_{n}(z)$ are analytic and $\neq0$ in a region $\Omega$, and if $f_{n}(z)$ converges to $f(z)$, uniformly on every compact subset of $\Omega$, then $f(z)$ is either identically zero or never equal to zero in $\Omega$.
\end{originalstatement}
Suppose that $f(z)$ is not identically zero. The zeros of $f(z)$ are in any case isolated. For any point $z_0 \in \Omega$ there is therefore a number $r > 0$ such that $f(z)$ is defined and $\neq 0$ for $0 < |z - z_0| \leq r$. In particular, $|f(z)|$ has a positive minimum on the circle $|z - z_0| = r$, which we denote by $C$. It follows that $1 / f_n(z)$ converges uniformly to $1 / f(z)$ on $C$. Since it is also true that $f_n'(z) \to f'(z)$, uniformly on $C$, we may conclude that
\begin{equation*}
\lim _ {n \rightarrow \infty} \frac {1}{2 \pi i} \int_ {C} \frac {f _ {n} ^ {\prime} (z)}{f _ {n} (z)} d z = \frac {1}{2 \pi i} \int_ {C} \frac {f ^ {\prime} (z)}{f (z)} d z.
\end{equation*}
But the integrals on the left are all zero, for they give the number of roots of the equation $f_{n}(z) = 0$ inside of C. The integral on the right is therefore zero, and consequently $f(z_{0}) \neq 0$ by the same interpretation of the integral. Since $z_{0}$ was arbitrary, the theorem follows.

\subsection*{EXERCISES}
\originalnumber{1}{exercise:5.1.1.1.1}
1. Using Taylor's theorem applied to a branch of $\log(1+z/n)$, prove that
\begin{equation*}
\lim _ {n \rightarrow \infty} \left(1 + \frac {z}{n}\right) ^ {n} = e ^ {z}
\end{equation*}
uniformly on all compact sets.

\originalnumber{2}{exercise:5.1.1.1.2}
2. Show that the series
\begin{equation*}
\zeta (z) = \sum_ {n = 1} ^ {\infty} n ^ {- z}
\end{equation*}
converges for $\operatorname{Re} z > 1$, and represent its derivative in series form.

\originalnumber{3}{exercise:5.1.1.1.3}
3. Prove that
\begin{equation*}
(1 - 2 ^ {1 - z}) \zeta (z) = 1 ^ {- z} - 2 ^ {- z} + 3 ^ {- z} - \dots
\end{equation*}
and that the latter series represents an analytic function for $\operatorname{Re} z > 0$.


% Original PDF physical page 195
\sourcepage{195}
\originalnumber{4}{exercise:5.1.1.1.4}
4. As a generalization of Theorem~\hyperref[thm:5.2]{2}, prove that if the $f_{n}(z)$ have at most $m$ zeros in $\Omega$, then $f(z)$ is either identically zero or has at most $m$ zeros.

\originalnumber{5}{exercise:5.1.1.1.5}
5. Prove that
\begin{equation*}
\sum_ {n = 1} ^ {\infty} \frac {n z ^ {n}}{1 - z ^ {n}} = \sum_ {n = 1} ^ {\infty} \frac {z ^ {n}}{(1 - z ^ {n}) ^ {2}}
\end{equation*}
for $|z| < 1$. (Develop in a double series and reverse the order of summation.)

\setcounter{section}{1}\setcounter{subsection}{1}\subsection{The Taylor Series}\label{sec:5.1.2}
We show now that every analytic function can be developed in a convergent Taylor series\index{Taylor series|(}\index{Taylor series}. This is an almost immediate consequence of the finite Taylor development given in Chap.~\hyperref[chap:4]{4}, Sec.~\hyperref[sec:4.3.1]{3.1}, Theorem~\hyperref[thm:4.8]{8}, together with the corresponding representation of the remainder term. According to this theorem, if $f(z)$ is analytic in a region $\Omega$ containing $z_{0}$, we can write
\begin{equation*}
\begin{aligned}
f(z)={}&f(z_0)+\frac{f'(z_0)}{1!}(z-z_0)+\dots\\
&+\frac{f^{(n)}(z_0)}{n!}(z-z_0)^n+f_{n+1}(z)(z-z_0)^{n+1}
\end{aligned}
\end{equation*}
with
\begin{equation*}
f _ {n + 1} (z) = \frac {1}{2 \pi i} \int_ {C} \frac {f (\zeta) d \zeta}{(\zeta - z _ {0}) ^ {n + 1} (\zeta - z)}.
\end{equation*}
In the last formula C is any circle $|z - z_{0}| = \rho$ such that the closed disk $|z - z_{0}| \leq \rho$ is contained in $\Omega$.

If $M$ denotes the maximum of $|f(z)|$ on $C$, we obtain at once the estimate
\begin{equation*}
\left| f _ {n + 1} (z) \left(z - z _ {0}\right) ^ {n + 1} \right| \leq \frac {M \left| z - z _ {0} \right| ^ {n + 1}}{\rho^ {n} (\rho - | z - z _ {0} |)}
\end{equation*}
We conclude that the remainder term tends uniformly to zero in every disk $|z - z_{0}| \leq r < \rho$. On the other hand, $\rho$ can be chosen arbitrarily close to the shortest distance from $z_{0}$ to the boundary of $\Omega$. We have proved:

\begin{originalstatement}{Theorem}{3}\label{thm:5.3}
If $f(z)$ is analytic in the region $\Omega$, containing $z_0$, then the representation
\begin{equation*}
f (z) = f \left(z _ {0}\right) + \frac {f ^ {\prime} \left(z _ {0}\right)}{1 !} \left(z - z _ {0}\right) + \dots + \frac {f ^ {(n)} \left(z _ {0}\right)}{n !} \left(z - z _ {0}\right) ^ {n} + \dots
\end{equation*}
is valid in the largest open disk of center $z_{0}$ contained in $\Omega$.

\end{originalstatement}

% Original PDF physical page 196
\index{Weierstrass's theorem|)}\sourcepage{196}
The radius of convergence of the Taylor series\index{Taylor series} is thus at least equal to the shortest distance from $z_{0}$ to the boundary of $\Omega$. It may well be larger, but if it is there is no guarantee that the series still represents $f(z)$ at all points which are simultaneously in $\Omega$ and in the circle of convergence.

We recall that the developments
\begin{equation*}
\begin{array}{l} e ^ {z} = 1 + z + \frac {z ^ {2}}{2 !} + \dots + \frac {z ^ {n}}{n !} + \dots \\ \cos z = 1 - \frac {z ^ {2}}{2 !} + \frac {z ^ {4}}{4 !} - \frac {z ^ {6}}{6 !} + \dots \\ \sin z = z - \frac {z ^ {3}}{3 !} + \frac {z ^ {5}}{5 !} - \frac {z ^ {7}}{7 !} + \dots . \end{array}
\end{equation*}
served as definitions of the functions they represent. Of course, as we have remarked before, every convergent power series is its own Taylor series. We gave earlier a direct proof that power series can be differentiated term by term. This is also a direct consequence of Weierstrass's theorem.

If we want to represent a fractional power of z or $\log z$ through a power series, we must first of all choose a well-defined branch, and secondly we have to choose a center $z_{0} \neq 0$. It amounts to the same thing if we develop the function $(1 + z)^{\mu}$ or $\log(1 + z)$ about the origin, choosing the branch which is respectively equal to 1 or 0 at the origin. Since this branch is single-valued and analytic in $|z| < 1$, the radius of convergence is at least 1. It is elementary to compute the coefficients, and we obtain
\begin{equation*}
\begin{array}{c} (1 + z) ^ {\mu} = 1 + \mu z + \binom {\mu} {2} z ^ {2} + \dots + \binom {\mu} {n} z ^ {n} + \dots \\ \log (1 + z) = z - \frac {z ^ {2}}{2} + \frac {z ^ {3}}{3} - \frac {z ^ {4}}{4} + \frac {z ^ {5}}{5} - \dots . \end{array}
\end{equation*}
where the binomial coefficients are defined by
\begin{equation*}
\binom {\mu} {n} = \frac {\mu (\mu - 1) \cdot \cdot \cdot (\mu - n + 1)}{1 \cdot 2 \cdot \cdot \cdot n}.
\end{equation*}
If the logarithmic series had a radius of convergence greater than 1, then $\log(1+z)$ would be bounded for $|z|<1$. Since this is not the case, the radius of convergence must be exactly 1. Similarly, if the binomial series were convergent in a circle of radius >1, the function $(1+z)^{\mu}$ and all its derivatives would be bounded in $|z|<1$. Unless $\mu$ is a positive integer, one of the derivatives will be a negative power of $1+z$, and hence unbounded. Thus the radius of convergence is precisely 1 except in the trivial case in which the binomial series reduces to a polynomial.


% Original PDF physical page 197
\sourcepage{197}
The series developments of the cyclometric functions arc tan z and arc sin z are most easily obtained by consideration of the derived series. From the expansion
\begin{equation*}
\frac {1}{1 + z ^ {2}} = 1 - z ^ {2} + z ^ {4} - z ^ {6} + \dots
\end{equation*}
we obtain by integration
\begin{equation*}
\operatorname{arc} \tan z = z - \frac {z ^ {3}}{3} + \frac {z ^ {5}}{5} - \frac {z ^ {7}}{7} + \dots
\end{equation*}
where the branch is uniquely determined as
\begin{equation*}
\operatorname{arc} \tan z = \int_ {0} ^ {z} \frac {d z}{1 + z ^ {2}}
\end{equation*}
for any path inside the unit circle. For justification we can either rely on uniform convergence or apply Theorem~\hyperref[thm:5.1]{1}. The radius of convergence cannot be greater than that of the derived series, and hence it is exactly 1.

If $\sqrt{1 - z^2}$ is the branch with a positive real part, we have
\begin{equation*}
\frac {1}{\sqrt {1 - z ^ {2}}} = 1 + \frac {1}{2} z ^ {2} + \frac {1 \cdot 3}{2 \cdot 4} z ^ {4} + \frac {1 \cdot 3 \cdot 5}{2 \cdot 4 \cdot 6} z ^ {6} + \dots
\end{equation*}
for $|z| < 1$, and through integration we obtain
\begin{equation*}
\operatorname{arc} \sin z = z + \frac {1}{2} \frac {z ^ {3}}{3} + \frac {1 \cdot 3}{2 \cdot 4} \frac {z ^ {5}}{5} + \frac {1 \cdot 3 \cdot 5}{2 \cdot 4 \cdot 6} \frac {z ^ {7}}{7} + \dots .
\end{equation*}
The series represents the principal branch of arc sin z with a real part between $-\pi/2$ and $\pi/2$.

For combinations of elementary functions it is mostly not possible to find a general law for the coefficients. In order to find the first few coefficients we need not, however, calculate the successive derivatives. There are simple techniques which allow us to compute, with a reasonable amount of labor, all the coefficients that we are likely to need.

It is convenient to introduce the notation $[z^{n}]$ for any function which is analytic and has a zero of at least order n at the origin; less precisely, $[z^{n}]$ denotes a function which “contains the factor $z^{n}$.” With this notation any function which is analytic at the origin can be written in the form
\begin{equation*}
f (z) = a _ {0} + a _ {1} z + \dots + a _ {n} z ^ {n} + [ z ^ {n + 1} ],
\end{equation*}
where the coefficients are uniquely determined and equal to the Taylor coefficients of $f(z)$. Thus, in order to find the first n coefficients of the Taylor expansion, it is sufficient to determine a polynomial $P_{n}(z)$ such
% Original PDF physical page 198
\sourcepage{198}
that $f(z) - P_n(z)$ has a zero of at least order $n + 1$ at the origin. The degree of $P_n(z)$ does not matter; it is true in any case that the coefficients of $z^m, m \leq n$, are the Taylor coefficients of $f(z)$.

For instance, suppose that
\begin{equation*}
\begin{array}{l} f (z) = a _ {0} + a _ {1} z + a _ {2} z ^ {2} + \cdot \cdot \cdot + a _ {n} z ^ {n} + \cdot \cdot \cdot \\ g (z) = b _ {0} + b _ {1} z + b _ {2} z ^ {2} + \cdot \cdot \cdot + b _ {n} z ^ {n} + \cdot \cdot \cdot . \end{array}
\end{equation*}
With an abbreviated notation we write
\begin{equation*}
f (z) = P _ {n} (z) + [ z ^ {n + 1} ]; \quad g (z) = Q _ {n} (z) + [ z ^ {n + 1} ].
\end{equation*}
It is then clear that $f(z)g(z) = P_{n}(z)Q_{n}(z) + [z^{n+1}]$, and the coefficients of the terms of degree $\leq n$ in $P_{n}Q_{n}$ are the Taylor coefficients of the product $f(z)g(z)$. Explicitly we obtain
\begin{equation*}
\begin{aligned}
f(z)g(z)={}&a_0b_0+(a_0b_1+a_1b_0)z+\dots\\
&+(a_0b_n+a_1b_{n-1}+\dots+a_nb_0)z^n+\dots.
\end{aligned}
\end{equation*}
In deriving this expansion we have not even mentioned the question of convergence, but since the development is identical with the Taylor development of $f(z)g(z)$, it follows by Theorem~\hyperref[thm:5.3]{3} that the radius of convergence is at least equal to the smaller of the radii of convergence of the given series $f(z)$ and $g(z)$. In the practical computation of $P_{n}Q_{n}$ it is of course not necessary to determine the terms of degree higher than n.

In the case of a quotient $f(z)/g(z)$ the same method can be applied, provided that $g(0)=b_{0}\neq0$. By use of ordinary long division, continued until the remainder contains the factor $z^{n+1}$, we can determine a polynomial $R_{n}$ such that $P_{n}=Q_{n}R_{n}+[z^{n+1}]$. Then $f-R_{n}g=[z^{n+1}]$, and since $g(0)\neq0$ we find that $f/g=R_{n}+[z^{n+1}]$. The coefficients of $R_{n}$ are the Taylor coefficients of $f(z)/g(z)$. They can be determined explicitly in determinant form, but the expressions are too complicated to be of essential help.

It is also important that we know how to form the development of a composite function $f(g(z))$. In this case, if $g(z)$ is developed around $z_{0}$, the expansion of $f(w)$ must be in powers of $w - g(z_{0})$. To simplify, let us assume that $z_{0} = 0$ and $g(0) = 0$. We can then set
\begin{equation*}
f (w) = a _ {0} + a _ {1} w + \dots + a _ {n} w ^ {n} + \dots
\end{equation*}
and $g(z) = b_{1}z + b_{2}z^{2} + \cdots + b_{n}z^{n} + \cdots$. Using the same notations as before we write $f(w) = P_{n}(w) + [w^{n+1}]$ and $g(z) = Q_{n}(z) + [z^{n+1}]$ with $Q_{n}(0) = 0$. Substituting $w = g(z)$ we have to observe that
\begin{equation*}
P _ {n} (Q _ {n} + [ z ^ {n + 1} ]) = P _ {n} (Q _ {n} (z)) + [ z ^ {n + 1} ]
\end{equation*}
% Original PDF physical page 199
\sourcepage{199}
and that any expression of the form $[w^{n+1}]$ becomes a $[z^{n+1}]$. Thus we obtain $f(g(z)) = P_{n}(Q_{n}(z)) + [z^{n+1}]$, and the Taylor coefficients of $f(g(z))$ are the coefficients of $P_{n}(Q_{n}(z))$ for powers $\leq n$.

Finally, we must be able to expand the inverse function of an analytic function $w = g(z)$. Here we may suppose that $g(0) = 0$, and we are looking for the branch of the inverse function $z = g^{-1}(w)$ which is analytic in a neighborhood of the origin and vanishes for $w = 0$. For the existence of the inverse function it is necessary and sufficient that $g'(0) \neq 0$; hence we assume that
\begin{equation*}
g (z) = a _ {1} z + a _ {2} z ^ {2} + \dots = Q _ {n} (z) + [ z ^ {n + 1} ]
\end{equation*}
with $a_1 \neq 0$. Our problem is to determine a polynomial $P_n(w)$ such that $P_n(Q_n(z)) = z + [z^{n+1}]$. In fact, under the assumption $a_1 \neq 0$ the notations $[z^{n+1}]$ and $[w^{n+1}]$ are interchangeable, and from $z = P_n(Q_n(z)) + [z^{n+1}]$ we obtain $z = P_n(g(z) + [z^{n+1}]) + [z^{n+1}] = P_n(w) + [w^{n+1}]$. Hence $P_n(w)$ determines the coefficients of $g^{-1}(w)$.

In order to prove the existence of a polynomial $P_{n}$ we proceed by induction. Clearly, we can take $P_{1}(w) = w/a_{1}$. If $P_{n-1}$ is given, we set $P_{n} = P_{n-1} + b_{n}w^{n}$ and obtain
\begin{equation*}
\begin{aligned}
&P_n(Q_n(z))=P_{n-1}(Q_n(z))+b_na_1^nz^n+[z^{n+1}]\\
&\quad=P_{n-1}(Q_{n-1}(z)+a_nz^n)+b_na_1^nz^n+[z^{n+1}]\\
&\quad=P_{n-1}(Q_{n-1}(z))+P_{n-1}'(Q_{n-1}(z))a_nz^n\\
&\qquad\qquad+b_na_1^nz^n+[z^{n+1}].
\end{aligned}
\end{equation*}
In the last member the first two terms form a known polynomial of the form $z + c_{n}z^{n} + [z^{n+1}]$, and we have only to take $b_{n} = -c_{n}a_{1}^{-n}$.

For practical purposes the development of the inverse function is found by successive substitutions. To illustrate the method we determine the expansion of $\tan w$ from the series
\begin{equation*}
w = \operatorname{arc} \tan z = z - \frac {z ^ {3}}{3} + \frac {z ^ {5}}{5} - \dots .
\end{equation*}
If we want the development to include fifth powers, we write
\begin{equation*}
z = w + \frac {z ^ {3}}{3} - \frac {z ^ {5}}{5} + [ z ^ {7} ]
\end{equation*}
and substitute this expression in the terms to the right. With appropriate remainders we obtain
\begin{equation*}
\begin{aligned}
z&=w+\frac13\left(w+\frac{z^3}3+[w^5]\right)^3\\
&\qquad-\frac15(w+[w^3])^5+[w^7]\\
&=w+\frac13w^3+\frac13w^2z^3-\frac15w^5+[w^7]\\
&=w+\frac13w^3+\frac13w^2(w+[w^3])^3-\frac15w^5+[w^7]\\
&=w+\frac13w^3+\frac{2}{15}w^5+[w^7].
\end{aligned}
\end{equation*}
% Original PDF physical page 200
\sourcepage{200}
Thus the development of tan w begins with the terms
\begin{equation*}
\tan w = w + \frac {1}{3} w ^ {3} + \frac {2}{15} w ^ {5} + \dots .
\end{equation*}
\subsection*{EXERCISES}
\originalnumber{1}{exercise:5.1.2.2.1}
1. Develop $1/(1+z^{2})$ in powers of $z-a$, $a$ being a real number. Find the general coefficient and for $a=1$ reduce to simplest form.

\originalnumber{2}{exercise:5.1.2.2.2}
2. The Legendre polynomials\index{Legendre polynomial} are defined as the coefficients $P_{n}(\alpha)$ in the development
\begin{equation*}
(1 - 2 \alpha z + z ^ {2}) ^ {- \frac {1}{2}} = 1 + P _ {1} (\alpha) z + P _ {2} (\alpha) z ^ {2} + \dots .
\end{equation*}
Find $P_{1}, P_{2}, P_{3}$, and $P_{4}$.

\originalnumber{3}{exercise:5.1.2.2.3}
3. Develop $\log(\sin z/z)$ in powers of $z$ up to the term $z^{6}$.

\originalnumber{4}{exercise:5.1.2.2.4}
4. What is the coefficient of $z^{7}$ in the Taylor development of $\tan z$?

\originalnumber{5}{exercise:5.1.2.2.5}
5. The Fibonacci numbers\index{Fibonacci numbers} are defined by $c_{0}=0$, $c_{1}=1$,
\begin{equation*}
c _ {n} = c _ {n - 1} + c _ {n - 2}.
\end{equation*}
Show that the $c_{n}$ are Taylor coefficients of a rational function, and determine a closed expression for $c_{n}$.

\setcounter{section}{1}\setcounter{subsection}{2}\subsection{The Laurent Series}\index{Laurent series|(}\index{Laurent series}\label{sec:5.1.3}
A series of the form
\begin{equation*}
b _ {0} + b _ {1} z ^ {- 1} + b _ {2} z ^ {- 2} + \dots + b _ {n} z ^ {- n} + \dots\tag{1}\label{eq:5.1}
\end{equation*}
can be considered as an ordinary power series in the variable 1/z. It will therefore converge outside of some circle $|z| = R$, except in the extreme case $R = \infty$; the convergence is uniform in every region $|z| \geq \rho > R$, and hence the series represents an analytic function in the region $|z| > R$. If the series \hyperref[eq:5.1]{(1)} is combined with an ordinary power series, we get a more general series of the form
\begin{equation*}
\sum_ {n = - \infty} ^ {+ \infty} a _ {n} z ^ {n}.\tag{2}\label{eq:5.2}
\end{equation*}
It will be termed convergent only if the parts consisting of nonnegative powers and negative powers are separately convergent. Since the first part converges in a disk $|z| < R_{2}$ and the second series in a region $|z| > R_{1}$, there is a common region of convergence only if $R_{1} < R_{2}$, and \hyperref[eq:5.2]{(2)} represents an analytic function in the annulus $R_{1} < |z| < R_{2}$.

Conversely, we may start from an analytic function $f(z)$ whose region of definition contains an annulus $R_{1} < |z| < R_{2}$, or more generally an annulus $R_{1} < |z - a| < R_{2}$. We shall show that such a function can


% Original PDF physical page 201
\index{Taylor series|)}\sourcepage{201}
always be developed in a general power series of the form
\begin{equation*}
f (z) = \sum_ {n = - \infty} ^ {+ \infty} A _ {n} (z - a) ^ {n}.
\end{equation*}
The proof is extremely simple. All we have to show is that $f(z)$ can be written as a sum $f_{1}(z) + f_{2}(z)$ where $f_{1}(z)$ is analytic for $|z - a| < R_{2}$ and $f_{2}(z)$ is analytic for $|z - a| > R_{1}$ with a removable singularity at $\infty$. Under these circumstances $f_{1}(z)$ can be developed in nonnegative powers of $z - a$, and $f_{2}(z)$ can be developed in nonnegative powers of $1 / (z - a)$.

To find the representation $f(z) = f_{1}(z) + f_{2}(z)$ define $f_{1}(z)$ by
\begin{equation*}
f _ {1} (z) = \frac {1}{2 \pi i} \int_ {| \zeta - a | = r} \frac {f (\zeta) d \zeta}{\zeta - z}
\end{equation*}
for $|z - a| < r < R_2$ and $f_2(z)$ by
\begin{equation*}
f _ {2} (z) = - \frac {1}{2 \pi i} \int_ {| \zeta - a | = r} \frac {f (\zeta) d \zeta}{\zeta - z}
\end{equation*}
for $R_{1} < r < |z - a|$. In both integrals the value of $r$ is irrelevant as long as the inequality is fulfilled, for it is an immediate consequence of Cauchy's theorem that the value of the integral does not change with $r$ provided that the circle does not pass over the point $z$. For this reason $f_{1}(z)$ and $f_{2}(z)$ are uniquely defined and represent analytic functions in $|z - a| < R_{2}$ and $|z - a| > R_{1}$ respectively. Moreover, by Cauchy's integral theorem $f(z) = f_{1}(z) + f_{2}(z)$.

The Taylor development of $f_{1}(z)$ is
\begin{equation*}
f _ {1} (z) = \sum_ {n = 0} ^ {\infty} A _ {n} (z - a) ^ {n}
\end{equation*}
with
\begin{equation*}
A _ {n} = \frac {1}{2 \pi i} \int_ {| \zeta - a | = r} \frac {f (\zeta) d \zeta}{(\zeta - a) ^ {n + 1}}.\tag{3}\label{eq:5.3}
\end{equation*}
In order to find the development of $f_{2}(z)$ we perform the transformation $\zeta = a + 1/\zeta'$, $z = a + 1/z'$. This transformation carries $|\zeta - a| = r$ into $|\zeta'| = 1/r$ with negative orientation, and by simple calculations we obtain
\begin{equation*}
f _ {2} \left(a + \frac {1}{z ^ {\prime}}\right) = \frac {1}{2 \pi i} \int_ {| \zeta^ {\prime} | = \frac {1}{r}} \frac {z ^ {\prime}}{\zeta^ {\prime}} \frac {f \left(a + \frac {1}{\zeta^ {\prime}}\right) d \zeta^ {\prime}}{\zeta^ {\prime} - z ^ {\prime}} = \sum_ {n = 1} ^ {\infty} B _ {n} z ^ {\prime n}
\end{equation*}
% Original PDF physical page 202
\sourcepage{202}
with
\begin{equation*}
B _ {n} = \frac {1}{2 \pi i} \int_ {| \zeta^ {\prime} | = \frac {1}{r}} \frac {f \left(a + \frac {1}{\zeta^ {\prime}}\right) d \zeta^ {\prime}}{\zeta^ {\prime n + 1}} = \frac {1}{2 \pi i} \int_ {| \zeta - a | = r} f (\zeta) (\zeta - a) ^ {n - 1} d \zeta .
\end{equation*}
This formula shows that we can write
\begin{equation*}
f (z) = \sum_ {n = - \infty} ^ {+ \infty} A _ {n} (z - a) ^ {n}
\end{equation*}
where all the coefficients $A_{n}$ are determined by \hyperref[eq:5.3]{(3)}. Observe that the integral in \hyperref[eq:5.3]{(3)} is independent of r as long as $R_{1} < r < R_{2}$.

If $R_{1}=0$ the point a is an isolated singularity and $A_{-1}=B_{1}$ is the residue at a, for $f(z)-A_{-1}(z-a)^{-1}$ is the derivative of a single-valued function in $0<|z-a|<R_{2}$.

\subsection*{EXERCISES}
\originalnumber{1}{exercise:5.1.3.3.1}
1. Prove that the Laurent development is unique.

\originalnumber{2}{exercise:5.1.3.3.2}
2. Let $\Omega$ be a doubly connected region whose complement consists of the components $E_1, E_2$. Prove that every analytic function $f(z)$ in $\Omega$ can be written in the form $f_1(z) + f_2(z)$ where $f_1(z)$ is analytic outside of $E_1$ and $f_2(z)$ is analytic outside of $E_2$. (The precise proof requires a construction like the one in Chap.~\hyperref[chap:4]{4}, Sec.~\hyperref[sec:4.4.5]{4.5}.)

\originalnumber{3}{exercise:5.1.3.3.3}
3. The expression
\begin{equation*}
\{f, z \} = \frac {f ^ {\prime \prime \prime} (z)}{f ^ {\prime} (z)} - \frac {3}{2} \left(\frac {f ^ {\prime \prime} (z)}{f ^ {\prime} (z)}\right) ^ {2}
\end{equation*}
is called the Schwarzian derivative\index{Schwarzian derivative} of $f$. If $f$ has a multiple zero or pole, find the leading term in the Laurent development of $\{f, z\}$. Answer: If $f(z) = a(z - z_0)^m + \cdots$, then $\{f, z\} = \frac{1}{2}(1 - m^2)(z - z_0)^{-2} + \cdots$.

\originalnumber{4}{exercise:5.1.3.3.4}
4. Show that the Laurent development of $(e^{z}-1)^{-1}$ at the origin is of the form
\begin{equation*}
\frac {1}{z} - \frac {1}{2} + \sum_ {1} ^ {\infty} (- 1) ^ {k - 1} \frac {B _ {k}}{(2 k) !} z ^ {2 k - 1}
\end{equation*}
where the numbers $B_{k}$ are known as the Bernoulli\index{Bernoulli J@Bernoulli, J.} numbers. Calculate $B_{1}, B_{2}, B_{3}$. (By Sec.~\hyperref[sec:5.2.1]{2.1}, Ex. \hyperref[exercise:5.2.1.4.5]{5}, the $B_{k}$ are all positive.)

\originalnumber{5}{exercise:5.1.3.3.5}
5. Express the Taylor development of $\tan z$ and the Laurent development of $\cot z$ in terms of the Bernoulli numbers.


% Original PDF physical page 203
\index{Laurent series|)}\sourcepage{203}
\setcounter{section}{1}\section{Partial Fractions And Factorization}\index{Fraction, partial|(}\index{Partial fraction|(}\label{sec:5.2}
A rational function has two standard representations, one by partial fractions and the other by factorization of the numerator and the denominator. The present section is devoted to similar representations of arbitrary meromorphic functions.

\setcounter{section}{2}\setcounter{subsection}{0}\subsection{Partial Fractions}\label{sec:5.2.1}
If the function $f(z)$ is meromorphic in a region $\Omega$, there corresponds to each pole $b_{\nu}$ a singular part of $f(z)$ consisting of the part of the Laurent development which contains the negative powers of $z - b_{\nu}$; it reduces to a polynomial $P_{\nu}(1/(z - b_{\nu}))$. It is tempting to subtract all singular parts in order to obtain a representation
\begin{equation*}
f (z) = \sum_ {\nu} P _ {\nu} \left(\frac {1}{z - b _ {\nu}}\right) + g (z)\tag{4}\label{eq:5.4}
\end{equation*}
where $g(z)$ would be analytic in $\Omega$. However, the sum on the right-hand side is in general infinite, and there is no guarantee that the series will converge. Nevertheless, there are many cases in which the series converges, and what is more, it is frequently possible to determine $g(z)$ explicitly from general considerations. In such cases the result is very rewarding; we obtain a simple expansion which is likely to be very helpful.

If the series in \hyperref[eq:5.4]{(4)} does not converge, the method needs to be modified. It is clear that nothing essential is lost if we subtract an analytic function $p_{\nu}(z)$ from each singular part $P_{\nu}$. By judicious choice of the functions $p_{\nu}$ the series $\sum_{\nu}(P_{\nu}-p_{\nu})$ can be made convergent. It is even possible to take the $p_{\nu}(z)$ to be polynomials.

We shall not prove the most general theorem to this effect. In the case where $\Omega$ is the whole plane we shall, however, prove that every meromorphic function has a development in partial fractions and, moreover, that the singular parts can be described arbitrarily. The theorem and its generalization to arbitrary regions are due to Mittag-Leffler\index{Mittag-Leffler G@Mittag-Leffler, G.}.

\begin{originalstatement}{Theorem}{4}\label{thm:5.4}
Let $\{b_{\nu}\}$ be a sequence of complex numbers with $\lim_{\nu \to \infty} b_{\nu} = \infty$,
and let $P_{\nu}(\zeta)$ be polynomials without constant term. Then there are functions which are meromorphic in the whole plane with poles at the points $b_{\nu}$ and the corresponding singular parts $P_{\nu}(1/(z - b_{\nu}))$. Moreover, the most general meromorphic function of this kind can be written in the form
\begin{equation*}
f (z) = \sum_ {\nu} \left[ P _ {\nu} \left(\frac {1}{z - b _ {\nu}}\right) - p _ {\nu} (z) \right] + g (z)\tag{5}\label{eq:5.5}
\end{equation*}
% Original PDF physical page 204
\sourcepage{204}
where the $p_{\nu}(z)$ are suitably chosen polynomials and $g(z)$ is analytic in the whole plane.

\end{originalstatement}
We may suppose that no $b_{\nu}$ is zero. The function $P_{\nu}(1 / (z - b_{\nu}))$ is analytic for $|z| < |b_{\nu}|$ and can thus be expanded in a Taylor series about the origin. We choose for $p_{\nu}(z)$ a partial sum of this series, ending, say, with the term of degree $n_{\nu}$. The difference $P_{\nu} - p_{\nu}$ can be estimated by use of the explicit expression for the remainder given in Chap.~\hyperref[chap:4]{4}, Sec.~\hyperref[sec:4.3.1]{3.1}. If the maximum of $|P_{\nu}|$ for $|z| \leq |b_{\nu}| / 2$ is denoted by $M_{\nu}$, we obtain
\begin{equation*}
\left| P _ {\nu} \left(\frac {1}{z - b}\right) - p _ {\nu} (z) \right| \leq 2 M \left(\frac {2 | z |}{| b _ {\nu} |}\right) ^ {n _ {\nu} + 1}\tag{6}\label{eq:5.6}
\end{equation*}
for all $|z| \leq |b_{\nu}| / 4$. By this estimate it is clear that the series in the right-hand member of \hyperref[eq:5.5]{(5)} can be made absolutely convergent in the whole plane, except at the poles, by choosing the $n_{\nu}$ sufficiently large. For instance, if we choose $n_{\nu}$ so large that $2^{n_{\nu}} \geq M_{\nu}2^{\nu}$, the estimate \hyperref[eq:5.6]{(6)} will show that the general term is majorized by $2^{-\nu}$ for all sufficiently large $\nu$.

Moreover, the estimate holds uniformly in any closed disk $|z| \leq R$, so that the convergence is actually uniform in that disk provided we omit the terms with $|b_{\nu}| \leq R$. By Weierstrass's theorem the remaining series represents an analytic function in $|z| \leq R$, and it follows that the full series is meromorphic in the whole plane with the singular parts $P_{\nu}(1/(z - b_{\nu}))$. The rest of the theorem is trivial.

As a first example we consider the function $\pi^{2}/\sin^{2}\pi z$, which has double poles at the points z = n for integral n. The singular part at the origin is $1/z^{2}$, and since $\sin^{2}\pi(z - n) = \sin^{2}\pi z$, the singular part at z = n is $1/(z - n)^{2}$. The series
\begin{equation*}
\sum_ {n = - \infty} ^ {+ \infty} \frac {1}{(z - n) ^ {2}}\tag{7}\label{eq:5.7}
\end{equation*}
is convergent for $z \neq n$, as seen by comparison with the familiar series $\sum_{1}^{\infty} 1/n^{2}$. It is uniformly convergent on any compact set after omission of the terms which become infinite on the set. For this reason we can write
\begin{equation*}
{\frac {\pi^ {2}}{\sin^ {2} \pi z}} = \sum_ {n = - \infty} ^ {+ \infty} {\frac {1}{(z - n) ^ {2}}} + g (z)\tag{8}\label{eq:5.8}
\end{equation*}
where $g(z)$ is analytic in the whole plane. We contend that $g(z)$ is identically zero.

To prove this we observe that the function $\pi^{2}/\sin^{2}\pi z$ and the series \hyperref[eq:5.7]{(7)} are both periodic with the period 1. Therefore the function $g(z)$ has the same period. For $z = x + iy$ we have (Chap.~\hyperref[chap:2]{2}, Sec.~\hyperref[sec:2.3.2]{3.2}, Ex. \hyperref[exercise:2.3.2.5.4]{4})


% Original PDF physical page 205
\sourcepage{205}
\begin{equation*}
\left| \sin \pi z \right| ^ {2} = \cosh^ {2} \pi y - \cos^ {2} \pi x
\end{equation*}
and hence $\pi^2/\sin^2\pi z$ tends uniformly to 0 as $|y|\to\infty$. But it is easy to see that the function \hyperref[eq:5.7]{(7)} has the same property. Indeed, the convergence is uniform for $|y|\geq1$, say, and the limit for $|y|\to\infty$ can thus be obtained by taking the limit in each term. We conclude that $g(z)$ tends uniformly to 0 for $|y|\to\infty$. This is sufficient to infer that $|g(z)|$ is bounded in a period strip $0\leq x\leq1$, and because of the periodicity $|g(z)|$ will be bounded in the whole plane. By Liouville's theorem $g(z)$ must reduce to a constant, and since the limit is 0 the constant must vanish. We have thus proved the identity
\begin{equation*}
\frac {\pi^ {2}}{\sin^ {2} \pi z} = \sum_ {- \infty} ^ {\infty} \frac {1}{(z - n) ^ {2}}.\tag{9}\label{eq:5.9}
\end{equation*}
From this equation a related identity can be obtained by integration. The left-hand member is the derivative of $-\pi \cot \pi z$, and the terms on the right are derivatives of $-1/(z - n)$. The series with the general term $1/(z - n)$ diverges, and a partial sum of the Taylor series must be subtracted from all the terms with $n \neq 0$. As it happens it is sufficient to subtract the constant terms, for the series
\begin{equation*}
\sum_ {n \neq 0} \left(\frac {1}{z - n} + \frac {1}{n}\right) = \sum_ {n \neq 0} \frac {z}{n (z - n)}
\end{equation*}
is comparable with $\sum_{1}^{\infty}1/n^{2}$ and hence convergent. The convergence is uniform on every compact set, provided that we omit the terms which become infinite. For this reason termwise differentiation is permissible, and we obtain
\begin{equation*}
\pi \cot \pi z = \frac {1}{z} + \sum_ {n \neq 0} \left(\frac {1}{z - n} + \frac {1}{n}\right)\tag{10}\label{eq:5.10}
\end{equation*}
except for an additive constant. If the terms corresponding to n and -n are bracketed together, \hyperref[eq:5.10]{(10)} can be written in the equivalent forms
\begin{equation*}
\pi \cot \pi z = \lim _ {m \rightarrow \infty} \sum_ {n = - m} ^ {m} \frac {1}{z - n} = \frac {1}{z} + \sum_ {n = 1} ^ {\infty} \frac {2 z}{z ^ {2} - n ^ {2}}.\tag{11}\label{eq:5.11}
\end{equation*}
With this way of writing it becomes evident that both members of the equation are odd functions of z, and for this reason the integration constant must vanish. The equations \hyperref[eq:5.10]{(10)} and \hyperref[eq:5.11]{(11)} are thus correctly stated.

Let us now reverse the procedure and try to evaluate the analogous


% Original PDF physical page 206
\sourcepage{206}
sum
\begin{equation*}
\lim _ {m \rightarrow \infty} \sum_ {- m} ^ {m} \frac {(- 1) ^ {n}}{z - n} = \frac {1}{z} + \sum_ {1} ^ {\infty} (- 1) ^ {n} \frac {2 z}{z ^ {2} - n ^ {2}}\tag{12}\label{eq:5.12}
\end{equation*}
which evidently represents a meromorphic function. It is very natural to separate the odd and even terms and write
\begin{equation*}
\sum_ {- (2 k + 1)} ^ {2 k + 1} \frac {(- 1) ^ {n}}{z - n} = \sum_ {n = - k} ^ {k} \frac {1}{z - 2 n} - \sum_ {n = - k - 1} ^ {k} \frac {1}{z - 1 - 2 n}.
\end{equation*}
By comparison with \hyperref[eq:5.11]{(11)} we find that the limit is
\begin{equation*}
\frac {\pi}{2} \cot \frac {\pi z}{2} - \frac {\pi}{2} \cot \frac {\pi (z - 1)}{2} = \frac {\pi}{\sin \pi z},
\end{equation*}
and we have proved that
\begin{equation*}
\frac {\pi}{\sin \pi z} = \lim _ {m \rightarrow \infty} \sum_ {- m} ^ {m} (- 1) ^ {n} \frac {1}{z - n}.\tag{13}\label{eq:5.13}
\end{equation*}
\subsection*{EXERCISES}
\originalnumber{1}{exercise:5.2.1.4.1}
1. Comparing coefficients in the Laurent developments of $\cot \pi z$ and of its expression as a sum of partial fractions, find the values of
\begin{equation*}
\sum_ {1} ^ {\infty} \frac {1}{n ^ {2}}, \quad \sum_ {1} ^ {\infty} \frac {1}{n ^ {4}}, \quad \sum_ {1} ^ {\infty} \frac {1}{n ^ {6}}.
\end{equation*}
Give a complete justification of the steps that are needed.

\originalnumber{2}{exercise:5.2.1.4.2}
2. Express
\begin{equation*}
\sum_ {- \infty} ^ {\infty} \frac {1}{z ^ {3} - n ^ {3}}
\end{equation*}
in closed form.

\originalnumber{3}{exercise:5.2.1.4.3}
3. Use \hyperref[eq:5.13]{(13)} to find the partial fraction\index{Fraction, partial}\index{Partial fraction} development of $1/\cos \pi z$, and show that it leads to $\pi/4 = 1 - \frac{1}{3} + \frac{1}{5} - \frac{1}{7} + \cdots$.

\originalnumber{4}{exercise:5.2.1.4.4}
4. What is the value of
\begin{equation*}
\sum_ {- \infty} ^ {\infty} \frac {1}{(z + n) ^ {2} + a ^ {2}}?
\end{equation*}
\originalnumber{5}{exercise:5.2.1.4.5}
5. Using the same method as in Ex.~\hyperref[exercise:5.2.1.4.1]{1}, show that
\begin{equation*}
\sum_ {1} ^ {\infty} \frac {1}{n ^ {2 k}} = 2 ^ {2 k - 1} \frac {B _ {k}}{(2 k) !} \pi^ {2 k}.
\end{equation*}
(See Sec.~\hyperref[sec:5.1.3]{1.3}, Ex. \hyperref[exercise:5.1.3.3.4]{4}, for the definition of $B_{k}$.)


% Original PDF physical page 207
\index{Fraction, partial|)}\index{Partial fraction|)}\sourcepage{207}
\setcounter{section}{2}\setcounter{subsection}{1}\subsection{Infinite Products}\label{sec:5.2.2}
An infinite product\index{Infinite product|(}\index{Infinite product} of complex numbers
\begin{equation*}
p _ {1} p _ {2} \cdot \cdot \cdot p _ {n} \cdot \cdot \cdot = \prod_ {n = 1} ^ {\infty} p _ {n}\tag{14}\label{eq:5.14}
\end{equation*}
is evaluated by taking the limit of the partial products $P_{n} = p_{1}p_{2} \cdots p_{n}$. It is said to converge to the value $P = \lim_{n \to \infty} P_{n}$ if this limit exists and is

different from zero. There are good reasons for excluding the value zero. For one thing, if the value P = 0 were permitted, any infinite product with one factor 0 would converge, and the convergence would not depend on the whole sequence of factors. On the other hand, in certain connections this convention is too radical. In fact, we wish to express a function as an infinite product, and this must be possible even if the function has zeros. For this reason we make the following agreement: The infinite product \hyperref[eq:5.14]{(14)} is said to converge if and only if at most a finite number of the factors are zero, and if the partial products formed by the nonvanishing factors tend to a finite limit which is different from zero.

In a convergent product the general factor $p_{n}$ tends to 1; this is clear by writing $p_{n} = P_{n}/P_{n-1}$, the zero factors being omitted. In view of this fact it is preferable to write all infinite products in the form
\begin{equation*}
\prod_ {n = 1} ^ {\infty} (1 + a _ {n})\tag{15}\label{eq:5.15}
\end{equation*}
so that $a_{n} \to 0$ is a necessary condition for convergence.

If no factor is zero, it is natural to compare the product \hyperref[eq:5.15]{(15)} with the infinite series
\begin{equation*}
\sum_ {n = 1} ^ {\infty} \log (1 + a _ {n}).\tag{16}\label{eq:5.16}
\end{equation*}
Since the $a_{n}$ are complex we must agree on a definite branch of the logarithms, and we decide to choose the principal branch in each term. Denote the partial sums of \hyperref[eq:5.16]{(16)} by $S_{n}$. Then $P_{n} = e^{S_{n}}$, and if $S_{n} \to S$ it follows that $P_{n}$ tends to the limit $P = e^{S}$ which is $\neq 0$. In other words, the convergence of \hyperref[eq:5.16]{(16)} is a sufficient condition for the convergence of \hyperref[eq:5.15]{(15)}.

In order to prove that the condition is also necessary, suppose that $P_{n} \rightarrow P \neq 0$. It is not true, in general, that the series \hyperref[eq:5.16]{(16)}, formed with the principal values, converges to the principal value of $\log P$; what we wish to show is that it converges to some value of $\log P$. For greater clarity we shall temporarily adopt the usage of denoting the principal value of the logarithm by Log and its imaginary part by Arg.


% Original PDF physical page 208
\sourcepage{208}
Because $P_{n} / P \to 1$ it is clear that $\operatorname{Log}(P_n / P) \to 0$ for $n \to \infty$. There exists an integer $h_n$ such that $\operatorname{Log}(P_n / P) = S_n - \operatorname{Log} P + h_n \cdot 2\pi i$. We pass to the differences to obtain $(h_{n+1} - h_n)2\pi i = \operatorname{Log}(P_{n+1} / P) - \operatorname{Log}(P_n / P) - \operatorname{Log}(1 + a_n)$ and hence $(h_{n+1} - h_n)2\pi = \operatorname{Arg}(P_{n+1} / P) - \operatorname{Arg}(P_n / P) - \operatorname{Arg}(1 + a_n)$. By definition, $|\operatorname{Arg}(1 + a_n)| \leq \pi$, and we know that $\operatorname{Arg}(P_{n+1} / P) - \operatorname{Arg}(P_n / P) \to 0$. For large $n$ this is incompatible with the previous equation unless $h_{n+1} = h_n$. Hence $h_n$ is ultimately equal to a fixed integer $h$, and it follows from $\operatorname{Log}(P_n / P) = S_n - \operatorname{Log} P + h \cdot 2\pi i$ that $S_n \to \operatorname{Log} P - h \cdot 2\pi i$. We have proved:

\begin{originalstatement}{Theorem}{5}\label{thm:5.5}
The infinite product\index{Infinite product} $\prod_{1}^{\infty}(1 + a_n)$ with $1 + a_{n} \neq 0$ converges simultaneously with the series $\sum_{1}^{\infty}\log (1 + a_n)$ whose terms represent the values of the principal branch of the logarithm.
\end{originalstatement}
The question of convergence of a product can thus be reduced to the more familiar question concerning the convergence of a series. It can be further reduced by observing that the series \hyperref[eq:5.16]{(16)} converges absolutely at the same time as the simpler series $\Sigma |a_{n}|$. This is an immediate consequence of the fact that
\begin{equation*}
\lim _ {z \rightarrow 0} \frac {\log (1 + z)}{z} = 1.
\end{equation*}
If either the series \hyperref[eq:5.16]{(16)} or $\sum_{1}^{\infty}|a_n|$ converges, we have $a_{n}\to 0$, and for a given $\varepsilon >0$ the double inequality
\begin{equation*}
(1 - \varepsilon) | a _ {n} | <   | \log (1 + a _ {n}) | <   (1 + \varepsilon) | a _ {n} |
\end{equation*}
will hold for all sufficiently large n. It follows immediately that the two series are in fact simultaneously absolutely convergent.

An infinite product is said to be absolutely convergent if and only if the corresponding series \hyperref[eq:5.16]{(16)} converges absolutely. With this terminology we can state our result in the following terms:

\begin{originalstatement}{Theorem}{6}\label{thm:5.6}
A necessary and sufficient condition for the absolute convergence of the product $\prod_{1}^{\infty}(1 + a_n)$ is the convergence of the series $\sum_{1}^{\infty}|a_n|$.
\end{originalstatement}
In the last theorem the emphasis is on absolute convergence. By


% Original PDF physical page 209
\sourcepage{209}
simple examples it can be shown that the convergence of $\sum_{1}^{\infty}a_n$ is neither sufficient nor necessary for the convergence of the product $\prod_{1}^{\infty}(1 + a_n)$.

It is clear what to understand by a uniformly convergent infinite product\index{Infinite product} whose factors are functions of a variable. The presence of zeros may cause some slight difficulties which can usually be avoided by considering only sets on which at most a finite number of the factors can vanish. If these factors are omitted, it is sufficient to study the uniform convergence of the remaining product. Theorems~\hyperref[thm:5.5]{5} and \hyperref[thm:5.6]{6} have obvious counterparts for uniform convergence. If we examine the proofs, we find that all estimates can be made uniform, and the conclusions lead to uniform convergence, at least on compact sets.

\subsection*{EXERCISES}
\originalnumber{1}{exercise:5.2.2.5.1}
1. Show that
\begin{equation*}
\prod_ {n = 2} ^ {\infty} \left(1 - \frac {1}{n ^ {2}}\right) = \frac {1}{2}.
\end{equation*}
\originalnumber{2}{exercise:5.2.2.5.2}
2. Prove that for $|z| < 1$
\begin{equation*}
(1 + z) (1 + z ^ {2}) (1 + z ^ {4}) (1 + z ^ {8}) \dots = \frac {1}{1 - z}.
\end{equation*}
\originalnumber{3}{exercise:5.2.2.5.3}
3. Prove that
\begin{equation*}
\prod_ {1} ^ {\infty} \left(1 + \frac {z}{n}\right) e ^ {- z / n}
\end{equation*}
converges absolutely and uniformly on every compact set.

\originalnumber{4}{exercise:5.2.2.5.4}
4. Prove that the value of an absolutely convergent product does not change if the factors are reordered.

\originalnumber{5}{exercise:5.2.2.5.5}
5. Show that the function
\begin{equation*}
\theta (z) = \prod_ {1} ^ {\infty} (1 + h ^ {2 n - 1} e ^ {z}) (1 + h ^ {2 n - 1} e ^ {- z})
\end{equation*}
where $|h| < 1$ is analytic in the whole plane and satisfies the functional equation
\begin{equation*}
\theta (z + 2 \log h) = h ^ {- 1} e ^ {- z} \theta (z).
\end{equation*}
\setcounter{section}{2}\setcounter{subsection}{2}\subsection{Canonical Products}\index{Canonical product|(}\index{Canonical product}\label{sec:5.2.3}
A function which is analytic in the whole plane is said to be \emph{entire}, or \emph{integral}. The simplest entire functions\index{Entire function}\index{Function!entire} which are not polynomials are $e^{z}$, $\sin z$, and $\cos z$.

If $g(z)$ is an entire function, then $f(z) = e^{g(z)}$ is entire and $\neq 0$. Conversely, if $f(z)$ is any entire function which is never zero, let us show
% Original PDF physical page 210
\index{Infinite product|)}\sourcepage{210}
that $f(z)$ is of the form $e^{g(z)}$. To this end we observe that the function $f'(z)/f(z)$, being analytic in the whole plane, is the derivative of an entire function $g(z)$. From this fact we infer, by computation, that $f(z)e^{-g(z)}$ has the derivative zero, and hence $f(z)$ is a constant multiple of $e^{g(z)}$; the constant can be absorbed in $g(z)$.

By this method we can also find the most general entire function with a finite number of zeros. Assume that $f(z)$ has m zeros at the origin (m may be zero), and denote the other zeros by $a_{1}, a_{2}, \ldots, a_{N}$, multiple zeros being repeated. It is then plain that we can write
\begin{equation*}
f (z) = z ^ {m} e ^ {g (z)} \prod_ {1} ^ {N} \left(1 - \frac {z}{a _ {n}}\right).
\end{equation*}
If there are infinitely many zeros, we can try to obtain a similar representation by means of an infinite product. The obvious generalization would be
\begin{equation*}
f (z) = z ^ {m} e ^ {g (z)} \prod_ {1} ^ {\infty} \left(1 - \frac {z}{a _ {n}}\right).\tag{17}\label{eq:5.17}
\end{equation*}
This representation is valid if the infinite product converges uniformly on every compact set. In fact, if this is so the product represents an entire function with zeros at the same points (except for the origin) and with the same multiplicities as $f(z)$. It follows that the quotient can be written in the form $z^{m}e^{g(z)}$.

The product in \hyperref[eq:5.17]{(17)} converges absolutely if and only if $\sum_{1}^{\infty} 1 / |a_n|$ is convergent, and in this case the convergence is also uniform in every closed disk $|z| \leq R$. It is only under this special condition that we can obtain a representation of the form \hyperref[eq:5.17]{(17)}.

In the general case convergence-producing factors must be introduced. We consider an arbitrary sequence of complex numbers $a_{n} \neq 0$ with $\lim_{n \to \infty} a_{n} = \infty$, and prove the existence of polynomials $p_{n}(z)$ such that
\begin{equation*}
\prod_ {1} ^ {\infty} \left(1 - \frac {z}{a _ {n}}\right) e ^ {p _ {n} (z)}\tag{18}\label{eq:5.18}
\end{equation*}
converges to an entire function. The product converges together with the series with the general term
\begin{equation*}
r _ {n} (z) = \log \left(1 - \frac {z}{a _ {n}}\right) + p _ {n} (z)
\end{equation*}
% Original PDF physical page 211
\sourcepage{211}
where the branch of the logarithm shall be chosen so that the imaginary part of $r_{n}(z)$ lies between $-\pi$ and $\pi$ (inclusive).

For a given R we consider only the terms with $|a_{n}| > R$. In the disk $|z| \leq R$ the principal branch of $\log(1 - z/a_{n})$ can be developed in a Taylor series
\begin{equation*}
\log \left(1 - \frac {z}{a _ {n}}\right) = - \frac {z}{a _ {n}} - \frac {1}{2} \left(\frac {z}{a _ {n}}\right) ^ {2} - \frac {1}{3} \left(\frac {z}{a _ {n}}\right) ^ {3} - \dots .
\end{equation*}
We reverse the signs and choose $p_n(z)$ as a partial sum
\begin{equation*}
p _ {n} (z) = \frac {z}{a _ {n}} + \frac {1}{2} \left(\frac {z}{a _ {n}}\right) ^ {2} + \dots + \frac {1}{m _ {n}} \left(\frac {z}{a _ {n}}\right) ^ {m _ {n}}.
\end{equation*}
Then $r_n(z)$ has the representation
\begin{equation*}
r _ {n} (z) = - \frac {1}{m _ {n} + 1} \left(\frac {z}{a _ {n}}\right) ^ {m _ {n} + 1} - \frac {1}{m _ {n} + 2} \left(\frac {z}{a _ {n}}\right) ^ {m _ {n} + 2} - \dots
\end{equation*}
and we obtain easily the estimate
\begin{equation*}
\left| r _ {n} (z) \right| \leq \frac {1}{m _ {n} + 1} \left(\frac {R}{\left| a _ {n} \right|}\right) ^ {m _ {n} + 1} \left(1 - \frac {R}{\left| a _ {n} \right|}\right) ^ {- 1}.\tag{19}\label{eq:5.19}
\end{equation*}
Suppose now that the series
\begin{equation*}
\sum_ {n = 1} ^ {\infty} \frac {1}{m _ {n} + 1} \left(\frac {R}{| a _ {n} |}\right) ^ {m _ {n} + 1}\tag{20}\label{eq:5.20}
\end{equation*}
converges. By the estimate \hyperref[eq:5.19]{(19)} it follows first that $r_{n}(z) \to 0$, and hence $r_{n}(z)$ has an imaginary part between $-\pi$ and $\pi$ as soon as n is sufficiently large. Moreover, the comparison shows that the series $\Sigma r_{n}(z)$ is absolutely and uniformly convergent for $|z| \leq R$, and thus the product \hyperref[eq:5.18]{(18)} represents an analytic function in $|z| < R$. For the sake of the reasoning we had to exclude the values $|a_{n}| \leq R$, but it is clear that the uniform convergence of \hyperref[eq:5.18]{(18)} is not affected when the corresponding factors are again taken into account.

It remains only to show that the series \hyperref[eq:5.20]{(20)} can be made convergent for all R. But this is obvious, for if we take $m_{n} = n$ it is clear that \hyperref[eq:5.20]{(20)} has a majorant geometric series with ratio < 1 for any fixed value of R.

\begin{originalstatement}{Theorem}{7}\label{thm:5.7}
There exists an entire function with arbitrarily prescribed zeros $a_{n}$ provided that, in the case of infinitely many zeros, $a_{n} \to \infty$. Every entire function with these and no other zeros can be written in the form
\begin{equation*}
f (z) = z ^ {m} e ^ {g (z)} \prod_ {n = 1} ^ {\infty} \left(1 - \frac {z}{a _ {n}}\right) e ^ {\frac {z}{a _ {n}} + \frac {1}{2} \left(\frac {z}{a _ {n}}\right) ^ {2} + \dots + \frac {1}{m _ {n}} \left(\frac {z}{a _ {n}}\right) ^ {m _ {n}}}\tag{21}\label{eq:5.21}
\end{equation*}
where the product is taken over all $a_{n} \neq 0$, the $m_{n}$ are certain integers, and $g(z)$ is an entire function.

\end{originalstatement}

% Original PDF physical page 212
\sourcepage{212}
This theorem is due to Weierstrass. It has the following important corollary:

\begin{originalstatement}{Corollary}{}\label{cor:5.unnumbered-p212-b1}
Every function which is meromorphic in the whole plane is the quotient of two entire functions.
\end{originalstatement}
In fact, if $F(z)$ is meromorphic in the whole plane, we can find an entire function $g(z)$ with the poles of $F(z)$ for zeros. The product $F(z)g(z)$ is then an entire function $f(z)$, and we obtain $F(z) = f(z)/g(z)$.

The representation \hyperref[eq:5.21]{(21)} becomes considerably more interesting if it is possible to choose all the $m_{n}$ equal to each other. The preceding proof has shown that the product
\begin{equation*}
\prod_ {1} ^ {\infty} \left(1 - \frac {z}{a _ {n}}\right) e ^ {\frac {z}{a _ {n}} + \frac {1}{2} \left(\frac {z}{a _ {n}}\right) ^ {2} + \dots + \frac {1}{h} \left(\frac {z}{a _ {n}}\right) ^ {h}}\tag{22}\label{eq:5.22}
\end{equation*}
converges and represents an entire function provided that the series $\sum_{n=1}^{\infty}(R/|a_{n}|)^{h+1}/(h+1)$ converges for all R, that is to say provided that $\Sigma1/|a_{n}|^{h+1}<\infty$. Assume that h is the smallest integer for which this series converges; the expression \hyperref[eq:5.22]{(22)} is then called the \emph{canonical product} associated with the sequence $\{a_{n}\}$, and h is the \emph{genus} of the canonical product\index{Canonical product}.

Whenever possible we use the canonical product in the representation \hyperref[eq:5.21]{(21)}, which is thereby uniquely determined. If in this representation $g(z)$ reduces to a polynomial, the function $f(z)$ is said to be of finite genus\index{Genus}, and the genus of $f(z)$ is by definition equal to the degree of this polynomial or to the genus of the canonical product, whichever is the larger. For instance, an entire function of genus zero is of the form
\begin{equation*}
C z ^ {m} \prod_ {1} ^ {\infty} \left(1 - \frac {z}{a _ {n}}\right)
\end{equation*}
with $\Sigma 1 / |a_n| < \infty$. The canonical representation of an entire function of genus 1 is either of the form
\begin{equation*}
C z ^ {m} e ^ {\alpha z} \prod_ {1} ^ {\infty} \left(1 - \frac {z}{a _ {n}}\right) e ^ {z / a _ {n}}
\end{equation*}
with $\Sigma 1 / |a_n|^2 < \infty$, $\Sigma 1 / |a_n| = \infty$, or of the form
\begin{equation*}
C z ^ {m} e ^ {\alpha z} \prod_ {1} ^ {\infty} \left(1 - \frac {z}{a _ {n}}\right)
\end{equation*}
% Original PDF physical page 213
\sourcepage{213}
with $\Sigma 1 / |a_n| < \infty, \alpha \neq 0$.

As an application we consider the product representation of $\sin \pi z$. The zeros are the integers $z = \pm n$. Since $\Sigma1/n$ diverges and $\Sigma1/n^{2}$ converges, we must take h = 1 and obtain a representation of the form
\begin{equation*}
\sin \pi z = z e ^ {g (z)} \prod_ {n \neq 0} \left(1 - \frac {z}{n}\right) e ^ {z / n}.
\end{equation*}
In order to determine $g(z)$ we form the logarithmic derivatives on both sides. We find
\begin{equation*}
\pi \cot \pi z = \frac {1}{z} + g ^ {\prime} (z) + \sum_ {n \neq 0} \left(\frac {1}{z - n} + \frac {1}{n}\right)
\end{equation*}
where the procedure is easy to justify by uniform convergence on any compact set which does not contain the points $z = n$. By comparison with the previous formula \hyperref[eq:5.10]{(10)} we conclude that $g'(z) = 0$. Hence $g(z)$ is a constant, and since $\lim_{z \to 0} \sin \pi z / z = \pi$ we must have $e^{g(z)} = \pi$, and thus
\begin{equation*}
\sin \pi z = \pi z \prod_ {n \neq 0} \left(1 - \frac {z}{n}\right) e ^ {z / n}.\tag{23}\label{eq:5.23}
\end{equation*}
In this representation the factors corresponding to n and -n can be bracketed together, and we obtain the simple form
\begin{equation*}
\sin \pi z = \pi z \prod_ {1} ^ {\infty} \left(1 - \frac {z ^ {2}}{n ^ {2}}\right).\tag{24}\label{eq:5.24}
\end{equation*}
It follows from \hyperref[eq:5.23]{(23)} that $\sin \pi z$ is an entire function of genus 1.

\subsection*{EXERCISES}
\originalnumber{1}{exercise:5.2.3.6.1}
1. Suppose that $a_{n} \to \infty$ and that the $A_{n}$ are arbitrary complex numbers. Show that there exists an entire function $f(z)$ which satisfies $f(a_{n}) = A_{n}$.

Hint: Let $g(z)$ be a function with simple zeros at the $a_{n}$. Show that
\begin{equation*}
\sum_ {1} ^ {\infty} g (z) \frac {e ^ {\gamma_ {n} (z - a _ {n})}}{z - a _ {n}} \cdot \frac {A _ {n}}{g ^ {\prime} (a _ {n})}
\end{equation*}
converges for some choice of the numbers $\gamma_{n}$.

\originalnumber{2}{exercise:5.2.3.6.2}
2. Prove that
\begin{equation*}
\sin \pi (z + \alpha) = e ^ {\pi z \cot \pi \alpha} \prod_ {- \infty} ^ {\infty} \left(1 + \frac {z}{n + \alpha}\right) e ^ {- z / (n + \alpha)}
\end{equation*}
% Original PDF physical page 214
\index{Canonical product|)}\sourcepage{214}
whenever $\alpha$ is not an integer. Hint: Denote the factor in front of the canonical product by $g(z)$ and determine $g'(z)/g(z)$.

\originalnumber{3}{exercise:5.2.3.6.3}
3. What is the genus of cos $\sqrt{z}$?

\originalnumber{4}{exercise:5.2.3.6.4}
4. If $f(z)$ is of genus $h$, how large and how small can the genus of $f(z^2)$ be?

\originalnumber{5}{exercise:5.2.3.6.5}
5. Show that if $f(z)$ is of genus 0 or 1 with real zeros, and if $f(z)$ is real for real $z$, then all zeros of $f'(z)$ are real. Hint: Consider $\operatorname{Im} f'(z) / f(z)$.

\setcounter{section}{2}\setcounter{subsection}{3}\subsection{The Gamma Function}\index{Function!gamma|(}\index{Gamma function|(}\index{Function!gamma}\index{Gamma function}\label{sec:5.2.4}
The function $\sin \pi z$ has all the integers for zeros, and it is the simplest function with this property. We shall now introduce functions which have only the positive or only the negative integers for zeros. The simplest function with, for instance, the negative integers for zeros is the corresponding canonical product
\begin{equation*}
G (z) = \prod_ {1} ^ {\infty} \left(1 + \frac {z}{n}\right) e ^ {- z / n}.\tag{25}\label{eq:5.25}
\end{equation*}
It is evident that $G(-z)$ has then the positive integers for zeros, and by comparison with the product representation \hyperref[eq:5.23]{(23)} of $\sin \pi z$ we find at once
\begin{equation*}
z G (z) G (- z) = \frac {\sin \pi z}{\pi}.\tag{26}\label{eq:5.26}
\end{equation*}
Because of the manner in which $G(z)$ has been constructed, it is bound to have other simple properties. We observe that $G(z - 1)$ has the same zeros as $G(z)$, and in addition a zero at the origin. It is therefore clear that we can write
\begin{equation*}
G (z - 1) = z e ^ {\gamma (z)} G (z),
\end{equation*}
where $\gamma(z)$ is an entire function. In order to determine $\gamma(z)$ we take the logarithmic derivatives on both sides. This gives the equation
\begin{equation*}
\sum_ {n = 1} ^ {\infty} \left(\frac {1}{z - 1 + n} - \frac {1}{n}\right) = \frac {1}{z} + \gamma^ {\prime} (z) + \sum_ {n = 1} ^ {\infty} \left(\frac {1}{z + n} - \frac {1}{n}\right).\tag{27}\label{eq:5.27}
\end{equation*}
In the series to the left we can replace n by $n + 1$. By this change we obtain
\begin{equation*}
\begin{aligned}
\sum_{n=1}^\infty\left(\frac1{z-1+n}-\frac1n\right)
&=\frac1z-1+\sum_{n=1}^\infty\left(\frac1{z+n}-\frac1{n+1}\right)\\
&=\frac1z-1+\sum_{n=1}^\infty\left(\frac1{z+n}-\frac1n\right)\\
&\qquad+\sum_{n=1}^\infty\left(\frac1n-\frac1{n+1}\right).
\end{aligned}
\end{equation*}
The last series has the sum 1, and hence equation \hyperref[eq:5.27]{(27)} reduces to $\gamma'(z) = 0$.


% Original PDF physical page 215
\sourcepage{215}
Thus $\gamma(z)$ is a constant, which we denote by $\gamma$, and $G(z)$ has the reproductive property $G(z - 1) = e^{\gamma}zG(z)$. It is somewhat simpler to consider the function $H(z) = G(z)e^{\gamma z}$ which evidently satisfies the functional equation $H(z - 1) = zH(z)$.

The value of $\gamma$ is easily determined. Taking $z = 1$ we have
\begin{equation*}
1 = G (0) = e ^ {\gamma} G (1),
\end{equation*}
and hence
\begin{equation*}
e ^ {- \gamma} = \prod_ {n = 1} ^ {\infty} \left(1 + \frac {1}{n}\right) e ^ {- 1 / n}.
\end{equation*}
Here the $n$th partial product can be written in the form
\begin{equation*}
(n + 1) e ^ {- (1 + \frac {1}{2} + \frac {1}{3} + \dots + 1 / n)},
\end{equation*}
and we obtain
\begin{equation*}
\gamma = \lim _ {n \rightarrow \infty} \left(1 + \frac {1}{2} + \frac {1}{3} + \dots + \frac {1}{n} - \log n\right).
\end{equation*}
The constant $\gamma$ is called Euler\index{Euler L@Euler, L.}'s constant; its approximate value is .57722.

If $H(z)$ satisfies $H(z - 1) = zH(z)$, then $\Gamma(z) = 1 / [zH(z)]$ satisfies $\Gamma(z - 1) = \Gamma(z) / (z - 1)$, or
\begin{equation*}
\Gamma (z + 1) = z \Gamma (z).\tag{28}\label{eq:5.28}
\end{equation*}
This is found to be a more useful relation, and for this reason it has become customary to implement the restricted stock of elementary functions by inclusion of $\Gamma(z)$ under the name of \emph{Euler's gamma function}\index{Function!gamma}\index{Gamma function}.

Our definition leads to the explicit representation
\begin{equation*}
\Gamma (z) = \frac {e ^ {- \gamma z}}{z} \prod_ {n = 1} ^ {\infty} \left(1 + \frac {z}{n}\right) ^ {- 1} e ^ {z / n},\tag{29}\label{eq:5.29}
\end{equation*}
and the formula \hyperref[eq:5.26]{(26)} takes the form
\begin{equation*}
\Gamma (z) \Gamma (1 - z) = \frac {\pi}{\sin \pi z}.\tag{30}\label{eq:5.30}
\end{equation*}
We observe that $\Gamma(z)$ is a meromorphic function with poles at $z = 0$, $-1, -2, \ldots$ but without zeros.

We have $\Gamma(1) = 1$, and by the functional equation we find $\Gamma(2) = 1$, $\Gamma(3) = 1 \cdot 2$, $\Gamma(4) = 1 \cdot 2 \cdot 3$ and generally $\Gamma(n) = (n - 1)!$. The $\Gamma$-function can thus be considered as a generalization of the factorial. From \hyperref[eq:5.30]{(30)} we conclude that $\Gamma(\frac{1}{2}) = \sqrt{\pi}$.

Other properties are most easily found by considering the second


% Original PDF physical page 216
\sourcepage{216}
derivative of $\log \Gamma(z)$ for which we find, by \hyperref[eq:5.29]{(29)}, the very simple expression
\begin{equation*}
\frac {d}{d z} \left(\frac {\Gamma^ {\prime} (z)}{\Gamma (z)}\right) = \sum_ {n = 0} ^ {\infty} \frac {1}{(z + n) ^ {2}}.\tag{31}\label{eq:5.31}
\end{equation*}
For instance, it is plain that $\Gamma(z)\,\Gamma(z + \frac{1}{2})$ and $\Gamma(2z)$ have the same poles, and by use of \hyperref[eq:5.31]{(31)} we find indeed that
\begin{equation*}
\begin{aligned}
&\frac{d}{dz}\left(\frac{\Gamma'(z)}{\Gamma(z)}\right)
+\frac{d}{dz}\left(\frac{\Gamma'(z+\frac12)}{\Gamma(z+\frac12)}\right)\\
&\quad=\sum_{n=0}^\infty\frac1{(z+n)^2}
+\sum_{n=0}^\infty\frac1{(z+n+\frac12)^2}\\
&\quad=4\left[\sum_{n=0}^\infty\frac1{(2z+2n)^2}
+\sum_{n=0}^\infty\frac1{(2z+2n+1)^2}\right]\\
&\quad=4\sum_{m=0}^\infty\frac1{(2z+m)^2}
=2\frac{d}{dz}\left(\frac{\Gamma'(2z)}{\Gamma(2z)}\right).
\end{aligned}
\end{equation*}
By integration we obtain
\begin{equation*}
\Gamma (z) \Gamma (z + \frac {1}{2}) = e ^ {a z + b} \Gamma (2 z),
\end{equation*}
where the constants $a$ and $b$ have yet to be determined. Substituting $z = \frac{1}{2}$ and $z = 1$ we make use of the known values $\Gamma(\frac{1}{2}) = \sqrt{\pi}$, $\Gamma(1) = 1$, $\Gamma(1\frac{1}{2}) = \frac{1}{2}\Gamma(\frac{1}{2}) = \frac{1}{2}\sqrt{\pi}$, $\Gamma(2) = 1$ and are led to the relations
\begin{equation*}
\frac {1}{2} a + b = \frac {1}{2} \log \pi , \quad a + b = \frac {1}{2} \log \pi - \log 2.
\end{equation*}
It follows that
\begin{equation*}
a = - 2 \log 2 \quad \text { and } \quad b = \frac {1}{2} \log \pi + \log 2;
\end{equation*}
the final result is thus
\begin{equation*}
\sqrt {\pi} \Gamma (2 z) = 2 ^ {2 z - 1} \Gamma (z) \Gamma (z + \frac {1}{2})
\end{equation*}
which is known as Legendre's duplication formula.

\subsection*{EXERCISES}
\originalnumber{1}{exercise:5.2.4.7.1}
1. Prove the formula of Gauss\index{Gauss K F@Gauss, K. F.}:
\begin{equation*}
(2 \pi) ^ {\frac {n - 1}{2}} \Gamma (z) = n ^ {z - \frac {1}{2}} \Gamma \left(\frac {z}{n}\right) \Gamma \left(\frac {z + 1}{n}\right) \cdot \cdot \cdot \Gamma \left(\frac {z + n - 1}{n}\right).
\end{equation*}
\originalnumber{2}{exercise:5.2.4.7.2}
2. Show that
\begin{equation*}
\Gamma \left(\frac {1}{6}\right) = 2 ^ {- \frac {1}{3}} \left(\frac {3}{\pi}\right) ^ {\frac {1}{2}} \Gamma \left(\frac {1}{3}\right) ^ {2}.
\end{equation*}
\originalnumber{3}{exercise:5.2.4.7.3}
3. What are the residues of $\Gamma(z)$ at the poles $z=-n$?


% Original PDF physical page 217
\sourcepage{217}
\setcounter{section}{2}\setcounter{subsection}{4}\subsection{Stirling's Formula}\index{Lindelöf E@Lindelöf, E.|(}\label{sec:5.2.5}
In most connections where the $\Gamma$ function can be applied, it is of utmost importance to have some information on the behavior of $\Gamma(z)$ for very large values of z. Fortunately, it is possible to calculate $\Gamma(z)$ with great precision and very little effort by means of a classical formula which goes under the name of Stirling's formula\index{Stirling's formula|(}\index{Stirling's formula}. There are many proofs of this formula. We choose to derive it by use of the residue calculus, following mainly the presentation of Lindelöf in his classical book on the calculus of residues. This is a very simple and above all a very instructive proof inasmuch as it gives us an opportunity to use residues in less trivial cases than previously.

The starting point is the formula \hyperref[eq:5.31]{(31)} for the second derivative of $\log \Gamma(z)$, and our immediate task is to express the partial sum
\begin{equation*}
\frac {1}{z ^ {2}} + \frac {1}{(z + 1) ^ {2}} + \frac {1}{(z + 2) ^ {2}} + \dots + \frac {1}{(z + n) ^ {2}}
\end{equation*}
as a convenient line integral. To this end we need a function with the residues $1/(z + \nu)^{2}$ at the integral points $\nu$; a good choice is
\begin{equation*}
\Phi (\zeta) = \frac {\pi \cot \pi \zeta}{(z + \zeta) ^ {2}}.
\end{equation*}
Here $\zeta$ is the variable while z enters only as a parameter, which in the first part of the derivation will be kept at a fixed value $z = x + iy$ with x > 0.

We apply the residue formula to the rectangle whose vertical sides lie on $\xi = 0$ and $\xi = n + \frac{1}{2}$ and with horizontal sides $\eta = \pm Y$, with the intention of letting first Y and then n tend to $\infty$. This contour, which we denote by K, passes through the pole at 0, but we know that the formula remains valid provided that we take the principal value of the integral and include one-half of the residue at the origin. Hence we obtain
\begin{equation*}
\operatorname{pr.v.} \frac {1}{2 \pi i} \int_ {K} \Phi (\zeta) d \zeta = - \frac {1}{2 z ^ {2}} + \sum_ {\nu = 0} ^ {n} \frac {1}{(z + \nu) ^ {2}}.
\end{equation*}
On the horizontal sides of the rectangle $\cot\pi\zeta$ tends uniformly to $\pm i$ for $Y\to\infty$. Since the factor $1/(z+\zeta)^{2}$ tends to zero, the corresponding integrals have the limit zero. We are now left with two integrals over infinite vertical lines. On each line $\xi=n+\frac{1}{2}$, $\cot\pi\xi$ is bounded, and because of the periodicity the bound is independent of n. The integral over the line $\xi=n+\frac{1}{2}$ is thus less than a constant times
\begin{equation*}
\int_ {\xi = n + \frac {1}{2}} \frac {d \eta}{| \zeta + z | ^ {2}}
\end{equation*}
% Original PDF physical page 218
\index{Lindelöf E@Lindelöf, E.|)}\sourcepage{218}
This integral can be evaluated, for on the line of integration
\begin{equation*}
\bar {\zeta} = 2 n + 1 - \zeta ,
\end{equation*}
and we obtain by residues
\begin{equation*}
\frac {1}{i} \int \frac {d \zeta}{| \zeta + z | ^ {2}} = \frac {1}{i} \int \frac {d \zeta}{(\zeta + z) (2 n + 1 - \zeta + \bar {z})} = \frac {2 \pi}{2 n + 1 + 2 x}.
\end{equation*}
The limit for $n \to \infty$ is thus zero.

Finally, the principal value of the integral over the imaginary axis from $-i\infty$ to $+i\infty$ can be written in the form
\begin{equation*}
\begin{aligned}
&\frac12\int_0^\infty\cot\pi i\eta
\left[\frac1{(i\eta+z)^2}-\frac1{(i\eta-z)^2}\right]d\eta\\
&\qquad=-\int_0^\infty\coth\pi\eta\cdot
\frac{2\eta z}{(\eta^2+z^2)^2}\,d\eta.
\end{aligned}
\end{equation*}
The sign has to be reversed, and we obtain the formula
\begin{equation*}
\frac {d}{d z} \left(\frac {\Gamma^ {\prime} (z)}{\Gamma (z)}\right) = \frac {1}{2 z ^ {2}} + \int_ {0} ^ {\infty} \coth \pi \eta \cdot \frac {2 \eta z}{(\eta^ {2} + z ^ {2}) ^ {2}} d \eta .\tag{32}\label{eq:5.32}
\end{equation*}
It is preferable to write
\begin{equation*}
\coth \pi \eta = 1 + \frac {2}{e ^ {2 \pi \eta} - 1}
\end{equation*}
and observe that the integral obtained from the term 1 has the value $1/z$. We can thus rewrite \hyperref[eq:5.32]{(32)} in the form
\begin{equation*}
\frac {d}{d z} \left(\frac {\Gamma^ {\prime} (z)}{\Gamma (z)}\right) = \frac {1}{z} + \frac {1}{2 z ^ {2}} + \int_ {0} ^ {\infty} \frac {4 \eta z}{(\eta^ {2} + z ^ {2}) ^ {2}} \cdot \frac {d \eta}{e ^ {2 \pi \eta} - 1}\tag{33}\label{eq:5.33}
\end{equation*}
where the integral is now very strongly convergent.

For z restricted to the right half plane this formula can be integrated. We find
\begin{equation*}
\frac {\Gamma^ {\prime} (z)}{\Gamma (z)} = C + \log z - \frac {1}{2 z} - \int_ {0} ^ {\infty} \frac {2 \eta}{\eta^ {2} + z ^ {2}} \cdot \frac {d \eta}{e ^ {2 \pi \eta} - 1},\tag{34}\label{eq:5.34}
\end{equation*}
where $\log z$ is the principal branch and C is an integration constant. The integration of the last term needs some justification. We have to make sure that the integral in \hyperref[eq:5.34]{(34)} can be differentiated under the sign of integration; this is so because the integral in \hyperref[eq:5.33]{(33)} converges uniformly when z is restricted to any compact set in the half plane x > 0.

We wish to integrate \hyperref[eq:5.34]{(34)} once more. This would obviously introduce arc tan $(z/\eta)$ in the integral, and although a single-valued branch could be defined we prefer to avoid the use of multiple-valued functions. That is possible if we first transform the integral in \hyperref[eq:5.34]{(34)} by partial integration. We obtain


% Original PDF physical page 219
\sourcepage{219}
\begin{equation*}
\int_ {0} ^ {\infty} \frac {2 \eta}{\eta^ {2} + z ^ {2}} \cdot \frac {d \eta}{e ^ {2 \pi \eta} - 1} = \frac {1}{\pi} \int_ {0} ^ {\infty} \frac {z ^ {2} - \eta^ {2}}{(\eta^ {2} + z ^ {2}) ^ {2}} \log (1 - e ^ {- 2 \pi \eta}) d \eta
\end{equation*}
where the logarithm is of course real. Now we can integrate with respect to $z$ and obtain
\begin{equation*}
\log \Gamma (z) = C ^ {\prime} + C z + \left(z - \frac {1}{2}\right) \log z + \frac {1}{\pi} \int_ {0} ^ {\infty} \frac {z}{\eta^ {2} + z ^ {2}} \log \frac {1}{1 - e ^ {- 2 \pi \eta}} d \eta \tag{35}\label{eq:5.35}
\end{equation*}
where $C'$ is a new integration constant and for convenience C - 1 has been replaced by C. The formula means that there exists, in the right half plane, a single-valued branch of $\log \Gamma(z)$ whose value is given by the right-hand member of the equation. By proper choice of $C'$ we obtain the branch of $\log \Gamma(z)$ which is real for real z.

It remains to determine the constants C and $C'$. To this end we must first study the behavior of the integral in \hyperref[eq:5.35]{(35)} which we denote by
\begin{equation*}
J (z) = \frac {1}{\pi} \int_ {0} ^ {\infty} \frac {z}{\eta^ {2} + z ^ {2}} \log \frac {1}{1 - e ^ {- 2 \pi \eta}} d \eta .\tag{36}\label{eq:5.36}
\end{equation*}
It is practically evident that $J(z) \to 0$ for $z \to \infty$ provided that z keeps away from the imaginary axis. Suppose for instance that z is restricted to the half plane $x \geq c > 0$. Breaking the integral into two parts we write
\begin{equation*}
J (z) = \int_ {0} ^ {\frac {| z |}{2}} + \int_ {\frac {| z |}{2}} ^ {\infty} = J _ {1} + J _ {2}.
\end{equation*}
In the first integral $|\eta^2 + z^2| \geq |z|^2 - |z/2|^2 = 3|z|^2/4$, and hence
\begin{equation*}
\left| J _ {1} \right| \leq \frac {4}{3 \pi | z |} \int_ {0} ^ {\infty} \log \frac {1}{1 - e ^ {- 2 \pi \eta}} d \eta .
\end{equation*}
In the second integral $|\eta^2 + z^2| = |z - i\eta| \cdot |z + i\eta| > c|z|$, and we find
\begin{equation*}
\left| J _ {2} \right| <   \frac {1}{\pi c} \int_ {\frac {| z |}{2}} ^ {\infty} \log \frac {1}{1 - e ^ {- 2 \pi \eta}} d \eta .
\end{equation*}
Since the integral of $\log\left(1-e^{-2\pi\eta}\right)$ is obviously convergent, we conclude that $J_{1}$ and $J_{2}$ tend to 0 as $z\to\infty$.

The value of C is found by substituting \hyperref[eq:5.35]{(35)} in the functional equation $\Gamma(z+1)=z\Gamma(z)$ or $\log\Gamma(z+1)=\log z+\log\Gamma(z)$; if we restrict z to positive values, there is no hesitancy about the branch of the logarithm. The substitution yields


% Original PDF physical page 220
\sourcepage{220}
\begin{equation*}
\begin{aligned}
&C'+Cz+C+(z+\tfrac12)\log(z+1)+J(z+1)\\
&\qquad=C'+Cz+(z+\tfrac12)\log z+J(z),
\end{aligned}
\end{equation*}
and this reduces to
\begin{equation*}
C = - \left(z + \frac {1}{2}\right) \log \left(1 + \frac {1}{z}\right) + J (z) - J (z + 1).
\end{equation*}
Letting $z \to \infty$ we find that $C = -1$.

Next we apply \hyperref[eq:5.35]{(35)} to the equation $\Gamma(z)\Gamma(1 - z) = \pi/\sin\pi z$, choosing $z = \frac{1}{2} + iy$. We obtain
\begin{equation*}
\begin{array}{r} 2 C ^ {\prime} -. 1 + i y \log (\frac {1}{2} + i y) - i y \log (\frac {1}{2} - i y) + J (\frac {1}{2} + i y) + J (\frac {1}{2} - i y) \\ = \log \pi - \log \cosh \pi y. \end{array}
\end{equation*}
This equation, in which the logarithms are to have their principal values, is so far proved only up to a constant multiple of $2\pi i$. But for y = 0 the equation is correct as it stands because \hyperref[eq:5.35]{(35)} determines the real value of $\log \Gamma(\frac{1}{2})$; hence it holds for all y. As $y \to \infty$ we known that $J(\frac{1}{2} + iy)$ and $J(\frac{1}{2} - iy)$ tend to 0. Developing the logarithm in a Taylor series we find
\begin{equation*}
i y \log \frac {\frac {1}{2} + i y}{\frac {1}{2} - i y} = i y \left(\pi i + \log \frac {1 + \frac {1}{2 i y}}{1 - \frac {1}{2 i y}}\right) = - \pi y + 1 + \varepsilon_ {1} (y)
\end{equation*}
while in the right-hand member
\begin{equation*}
\log \cosh \pi y = \pi y - \log 2 + \varepsilon_ {2} (y)
\end{equation*}
with $\varepsilon_1(y)$ and $\varepsilon_2(y)$ tending to 0. These considerations yield the value $C' = \frac{1}{2} \log 2\pi$. We have thus proved Stirling's formula\index{Stirling's formula} in the form
\begin{equation*}
\log \Gamma (z) = \frac {1}{2} \log 2 \pi - z + (z - \frac {1}{2}) \log z + J (z)\tag{37}\label{eq:5.37}
\end{equation*}
or equivalently
\begin{equation*}
\Gamma (z) = \sqrt {2 \pi} z ^ {z - \frac {1}{2}} e ^ {- z} e ^ {J (z)}\tag{38}\label{eq:5.38}
\end{equation*}
with the representation \hyperref[eq:5.36]{(36)} of the remainder valid in the right half plane. We know that $J(z)$ tends to 0 when $z \to \infty$ in a half plane $x \geq c > 0$.

In the expression for $J(z)$ we can develop the integrand in powers of 1/z and obtain
\begin{equation*}
J (z) = \frac {C _ {1}}{z} + \frac {C _ {2}}{z ^ {3}} + \dots + \frac {C _ {k}}{z ^ {2 k - 1}} + J _ {k} (z)
\end{equation*}
with


% Original PDF physical page 221
\sourcepage{221}
\begin{equation*}
C _ {\nu} = (- 1) ^ {\nu - 1} \frac {1}{\pi} \int_ {0} ^ {\infty} \eta^ {2 \nu - 2} \log \frac {1}{1 - e ^ {- 2 \pi \eta}} d \eta\tag{39}\label{eq:5.39}
\end{equation*}
and
\begin{equation*}
J _ {k} (z) = \frac {(- 1) ^ {k}}{z ^ {2 k + 1}} \frac {1}{\pi} \int_ {0} ^ {\infty} \frac {\eta^ {2 k}}{1 + (\eta / z) ^ {2}} \log \frac {1}{1 - e ^ {- 2 \pi \eta}} d \eta .
\end{equation*}
It can be proved (for instance by means of residues) that the coefficients $C_{\nu}$ are connected with the Bernoulli\index{Bernoulli J@Bernoulli, J.} numbers (cf. Ex.~\hyperref[exercise:5.1.3.3.4]{4}, Sec.~\hyperref[sec:5.1.3]{1.3}) by
\begin{equation*}
C _ {\nu} = (- 1) ^ {\nu - 1} \frac {1}{(2 \nu - 1) 2 \nu} B _ {\nu}.\tag{40}\label{eq:5.40}
\end{equation*}
Thus the development of $J(z)$ takes the form
\begin{equation*}
\begin{array}{r l} J (z) & = \frac {B _ {1}}{1 \cdot 2} \frac {1}{z} - \frac {B _ {2}}{3 \cdot 4} \cdot \frac {1}{z ^ {3}} + \dots \\ & \quad + (- 1) ^ {k - 1} \frac {B _ {k}}{(2 k - 1) 2 k} \frac {1}{z ^ {2 k - 1}} + J _ {k} (z). \end{array}\tag{41}\label{eq:5.41}
\end{equation*}
The reader is warned not to confuse this with a Laurent development. The function $J(z)$ is not defined in a neighborhood of $\infty$ and, therefore, does not have a Laurent development; moreover, if $k \to \infty$, the series obtained from \hyperref[eq:5.41]{(41)} does not converge. What we can say is that for a fixed k the expression $J_{k}(z)z^{2k}$ tends to 0 for $z \to \infty$ (in a half plane $x \geq c > 0$ ). This fact characterizes \hyperref[eq:5.41]{(41)} as an asymptotic development\index{Asymptotic development}. Such developments are very valuable when z is large in comparison with k, but for fixed z there is no advantage in letting k become very large.

Stirling's formula\index{Stirling's formula} can be used to prove that
\begin{equation*}
\Gamma (z) = \int_ {0} ^ {\infty} e ^ {- t} t ^ {z - 1} d t\tag{42}\label{eq:5.42}
\end{equation*}
whenever the integral converges, that is to say for x > 0. Until the identity has been proved, let the integral in \hyperref[eq:5.42]{(42)} be denoted by $F(z)$. Integrating by parts we find at once that
\begin{equation*}
F (z + 1) = \int_ {0} ^ {\infty} e ^ {- t} t ^ {z} d t = z \int_ {0} ^ {\infty} e ^ {- t} t ^ {z - 1} d t = z F (z).
\end{equation*}
Hence $F(z)$ satisfies the same functional equation as $\Gamma(z)$, and we find that $F(z)/\Gamma(z)=F(z+1)/\Gamma(z+1)$. In other words $F(z)/\Gamma(z)$ is periodic with the period 1. This shows, incidentally, that $F(z)$ can be defined in the whole plane although the integral representation is valid only in a half plane.

In order to prove that $F(z) / \Gamma(z)$ is constant we have to estimate $|F / \Gamma|$ in a period strip, for instance in the strip $1 \leq x \leq 2$. In the first place we have by \hyperref[eq:5.42]{(42)}


% Original PDF physical page 222
\index{Function!gamma|)}\sourcepage{222}
\begin{equation*}
| F (z) | \leq \int_ {0} ^ {\infty} e ^ {- t} t ^ {x - 1} d t = F (x),
\end{equation*}\index{Hadamard's theorem|(}
and hence $F(z)$ is bounded in the strip. Next, we use Stirling's formula\index{Stirling's formula} to find a lower bound of $|\Gamma(z)|$ for large $y$. From \hyperref[eq:5.37]{(37)} we obtain
\begin{equation*}
\log | \Gamma (z) | = \frac {1}{2} \log 2 \pi - x + (x - \frac {1}{2}) \log | z | - y \arg z + \operatorname{Re} J (z).
\end{equation*}
Only the term $-y$ arg $z$ becomes negatively infinite, being comparable to $-\pi |y| / 2$. Thus $|F / \Gamma|$ does not grow much more rapidly than $e^{\pi |y| / 2}$.

For an arbitrary function this would not suffice to conclude that the function must be constant, but for a function of period 1 it is more than enough. In fact, it is clear that $F / \Gamma$ can be expressed as a single-valued function of the variable $\zeta = e^{2\pi iz}$; to every value of $\zeta \neq 0$ there correspond infinitely many values of $z$ which differ by multiples of 1, and thus a single value of $F / \Gamma$. The function has isolated singularities at $\zeta = 0$ and $\zeta = \infty$, and our estimate shows that $|F / \Gamma|$ grows at most like $|\zeta|^{-\frac{1}{2}}$ for $\zeta \to 0$ and $|\zeta|^{\frac{1}{2}}$ for $\zeta \to \infty$. It follows that both singularities are removable, and hence $F / \Gamma$ must reduce to a constant. Finally, the fact that $F(1) = \Gamma(1) = 1$ shows that $F(z) = \Gamma(z)$.

\subsection*{EXERCISES}
\originalnumber{1}{exercise:5.2.5.8.1}
1. Prove the development \hyperref[eq:5.41]{(41)}.

\originalnumber{2}{exercise:5.2.5.8.2}
2. For real $x > 0$ prove that
\begin{equation*}
\Gamma (x) = \sqrt {2 \pi} x ^ {x - \frac {1}{2}} e ^ {- x} e ^ {\theta (x) / 12 x}
\end{equation*}
with $0 < \theta(x) < 1$.

\originalnumber{3}{exercise:5.2.5.8.3}
3. The formula \hyperref[eq:5.42]{(42)} permits us to evaluate the probability integral\index{Probability integral}
\begin{equation*}
\int_ {0} ^ {\infty} e ^ {- t ^ {2}} d t = \frac {1}{2} \int_ {0} ^ {\infty} e ^ {- x} x ^ {- \frac {1}{2}} d x = \frac {1}{2} \Gamma (\frac {1}{2}) = \frac {1}{2} \sqrt {\pi}.
\end{equation*}
Use this result together with Cauchy's theorem to compute the Fresnel integrals\index{Fresnel integral}
\begin{equation*}
\int_ {0} ^ {\infty} \sin (x ^ {2}) d x, \quad \int_ {0} ^ {\infty} \cos (x ^ {2}) d x.
\end{equation*}
Answer: Both are equal to $\frac{1}{2}\sqrt{\pi / 2}$.

\setcounter{section}{2}\section{Entire Functions}\label{sec:5.3}
In Sec.~\hyperref[sec:5.2.3]{2.3} we have already considered the representation of entire functions\index{Entire function|(}\index{Function!entire|(}\index{Entire function}\index{Function!entire} as infinite products, and, in special cases, as canonical products. In this section we study the connection between the product representation and the rate of growth of the function. Such questions were first investigated by Hadamard\index{Hadamard J@Hadamard, J.} who applied the results to his celebrated proof
% Original PDF physical page 223
\index{Gamma function|)}\index{Stirling's formula|)}\sourcepage{223}
of the Prime Number Theorem. Space does not permit us to include this application, but the basic importance of Hadamard's factorization theorem will be quite evident.

\setcounter{section}{3}\setcounter{subsection}{0}\subsection{Jensen's Formula}\index{Jensen's formula|(}\index{Jensen's formula}\label{sec:5.3.1}
If $f(z)$ is an analytic function, then $\log |f(z)|$ is harmonic except at the zeros of $f(z)$. Therefore, if $f(z)$ is analytic and free from zeros in $|z| \leq \rho$,
\begin{equation*}
\log | f (0) | = \frac {1}{2 \pi} \int_ {0} ^ {2 \pi} \log | f (\rho e ^ {i \theta}) | d \theta ,\tag{43}\label{eq:5.43}
\end{equation*}
and $\log |f(z)|$ can be expressed by Poisson's formula.

The equation \hyperref[eq:5.43]{(43)} remains valid if $f(z)$ has zeros on the circle $|z| = \rho$. The simplest proof is by dividing $f(z)$ with one factor $z - \rho e^{i\theta_{0}}$ for each zero. It is sufficient to show that
\begin{equation*}
\log \rho = \frac {1}{2 \pi} \int_ {0} ^ {2 \pi} \log | \rho e ^ {i \theta} - \rho e ^ {i \theta_ {0}} | d \theta
\end{equation*}
or
\begin{equation*}
\int_ {0} ^ {2 \pi} \log | e ^ {i \theta} - e ^ {i \theta_ {0}} | d \theta = 0.
\end{equation*}
This integral is evidently independent of $\theta_{0}$, and we have only to show that
\begin{equation*}
\int_ {0} ^ {2 \pi} \log | 1 - e ^ {i \theta} | d \theta = 0.
\end{equation*}
But this is a consequence of the formula
\begin{equation*}
\int_ {0} ^ {\pi} \log \sin x d x = - \pi \log 2
\end{equation*}
proved in Chap.~\hyperref[chap:4]{4}, Sec.~\hyperref[sec:4.5.3]{5.3} (cf. Chap.~\hyperref[chap:4]{4}, Sec.~\hyperref[sec:4.6.4]{6.4}, Ex. \hyperref[exercise:4.6.4.12.5]{5}).

We will now investigate what becomes of \hyperref[eq:5.43]{(43)} in the presence of zeros in the interior $|z| < \rho$. Denote these zeros by $a_{1}, a_{2}, \ldots, a_{n}$, multiple zeros being repeated, and assume first that z = 0 is not a zero. Then the function
\begin{equation*}
F (z) = f (z) \prod_ {i = 1} ^ {n} \frac {\rho^ {2} - \bar {a} _ {i} z}{\rho (z - a _ {i})}
\end{equation*}
is free from zeros in the disk, and $|F(z)| = |f(z)|$ on $|z| = \rho$. Consequently we obtain
\begin{equation*}
\log | F (0) | = \frac {1}{2 \pi} \int_ {0} ^ {2 \pi} \log | f (\rho e ^ {i \theta}) | d \theta
\end{equation*}
and, substituting the value of $F(0)$,


% Original PDF physical page 224
\sourcepage{224}
\begin{equation*}
\log | f (0) | = - \sum_ {i = 1} ^ {n} \log \left(\frac {\rho}{| a _ {i} |}\right) + \frac {1}{2 \pi} \int_ {0} ^ {2 \pi} \log | f (\rho e ^ {i \theta}) | d \theta .\tag{44}\label{eq:5.44}
\end{equation*}\index{Order, algebraic!of entire function|(}
This is known as Jensen's formula\index{Jensen's formula}. Its importance lies in the fact that it relates the modulus $|f(z)|$ on a circle to the moduli of the zeros.

If $f(0) = 0$, the formula is somewhat more complicated. Writing $f(z) = cz^h + \cdots$ we apply \hyperref[eq:5.44]{(44)} to $f(z)(\rho / z)^h$ and find that the left-hand member must be replaced by $\log |c| + h \log \rho$.

There is a similar generalization of Poisson's formula. All that is needed is to apply the ordinary Poisson formula to $\log |F(z)|$. We obtain
\begin{equation*}
\begin{aligned}
\log|f(z)|={}&-\sum_{i=1}^n\log\left|\frac{\rho^2-\bar a_i z}{\rho(z-a_i)}\right|\\
&+\frac1{2\pi}\int_0^{2\pi}\operatorname{Re}
\frac{\rho e^{i\theta}+z}{\rho e^{i\theta}-z}\log|f(\rho e^{i\theta})|\,d\theta,
\end{aligned}\tag{45}\label{eq:5.45}
\end{equation*}
provided that $f(z) \neq 0$. Equation \hyperref[eq:5.45]{(45)} is frequently referred to as the Poisson-Jensen formula\index{Poisson-Jensen formula}.

Strictly speaking the proof is valid only if $f \neq 0$ on $|z| = \rho$. But \hyperref[eq:5.44]{(44)} shows that the integral on the right is a continuous function of $\rho$, and from there it is easy to infer that the integral in \hyperref[eq:5.45]{(45)} is likewise continuous. In the general case \hyperref[eq:5.45]{(45)} can therefore be derived by letting $\rho$ approach a limit.

The Jensen and Poisson-Jensen formulas have important applications in the theory of entire functions\index{Entire function}\index{Function!entire}.

\setcounter{section}{3}\setcounter{subsection}{1}\subsection{Hadamard's Theorem}\index{Hadamard's theorem}\label{sec:5.3.2}
Let $f(z)$ be an entire function, and denote its zeros by $a_{n}$; for the sake of simplicity we will assume that $f(0) \neq 0$. We recall that the genus of an entire function (Sec.~\hyperref[sec:5.2.3]{2.3}) is the smallest integer $h$ such that $f(z)$ can be represented in the form
\begin{equation*}
f (z) = e ^ {g (z)} \prod_ {n} \left(1 - \frac {z}{a _ {n}}\right) e ^ {z / a _ {n} + \frac {1}{2} (z / a _ {n}) ^ {2} + \dots + (1 / h) (z / a _ {n}) ^ {h}}\tag{46}\label{eq:5.46}
\end{equation*}
where $g(z)$ is a polynomial of degree $\leq h$. If there is no such representation, the genus is infinite.

Denote by $M(r)$ the maximum of $|f(z)|$ on $|z| = r$. The order of the entire function $f(z)$ is defined by
\begin{equation*}
\lambda = \varlimsup_ {r \rightarrow \infty} \frac {\log \log M (r)}{\log r}.
\end{equation*}
According to this definition $\lambda$ is the smallest number such that
\begin{equation*}
M (r) \leq e ^ {r ^ {\lambda + \epsilon}}\tag{47}\label{eq:5.47}
\end{equation*}
for any given $\varepsilon > 0$ as soon as $r$ is sufficiently large.


% Original PDF physical page 225
\index{Jensen's formula|)}\index{Order, algebraic!of entire function|)}\sourcepage{225}
The genus and the order are closely related, as seen by the following theorem:

\begin{originalstatement}{Theorem}{8}\label{thm:5.8}
The genus and the order of an entire function\index{Entire function}\index{Function!entire} satisfy the double inequality $h \leq \lambda \leq h + 1$.
\end{originalstatement}
Assume first that $f(z)$ is of finite genus $h$. The exponential factor in \hyperref[eq:5.46]{(46)} is quite obviously of order $\leq h$, and the order of a product cannot exceed the orders of both factors. Hence it is sufficient to show that the canonical product is of order $\leq h + 1$. The convergence of the canonical product implies $\sum_{n} |a_n|^{-h-1} < \infty$; this is the essential hypothesis.

We denote the canonical product by $P(z)$ and write the individual factors as $E_{h}(z/a_{n})$ where
\begin{equation*}
E _ {h} (u) = (1 - u) e ^ {u + \frac {1}{2} u ^ {2} + \dots + (1 / h) u ^ {h}}
\end{equation*}
with the understanding that $E_0(u) = 1 - u$. We will show that
\begin{equation*}
\log | E _ {h} (u) | \leq (2 h + 1) | u | ^ {h + 1}\tag{48}\label{eq:5.48}
\end{equation*}
for all u.

If $|u| < 1$ we have by power-series development
\begin{equation*}
\log | E _ {h} (u) | \leq \frac {| u | ^ {h + 1}}{h + 1} + \frac {| u | ^ {h + 2}}{h + 2} + \dots \leq \frac {1}{h + 1} \frac {| u | ^ {h + 1}}{1 - | u |}
\end{equation*}
and thus
\begin{equation*}
(1 - | u |) \log | E _ {h} (u) | \leq | u | ^ {h + 1}.\tag{49}\label{eq:5.49}
\end{equation*}
For arbitrary $u$ and $h \geq 1$ it is also clear that
\begin{equation*}
\log | E _ {h} (u) | \leq \log | E _ {h - 1} (u) | + | u | ^ {h}.\tag{50}\label{eq:5.50}
\end{equation*}
The truth of \hyperref[eq:5.48]{(48)} is seen by induction. For h = 0 we need merely note that $\log |1 - u| \leq \log (1 + |u|) \leq |u|$. We assume \hyperref[eq:5.48]{(48)} with h - 1 in the place of h, that is to say
\begin{equation*}
\log | E _ {h - 1} (u) | \leq (2 h - 1) | u | ^ {h}.\tag{51}\label{eq:5.51}
\end{equation*}
It follows from \hyperref[eq:5.50]{(50)} and \hyperref[eq:5.51]{(51)} that $\log |E_{h}(u)| \leq 2h|u|^{h}$, and if $|u| \geq 1$, this implies \hyperref[eq:5.48]{(48)}. But if $|u| < 1$ we can also use \hyperref[eq:5.49]{(49)}, and together with \hyperref[eq:5.50]{(50)} and \hyperref[eq:5.51]{(51)} we obtain
\begin{equation*}
\log | E _ {h} (u) | \leq | u | \log | E _ {h - 1} (u) | + 2 | u | ^ {h + 1} \leq (2 h + 1) | u | ^ {h + 1}.
\end{equation*}
This completes the induction.


% Original PDF physical page 226
\sourcepage{226}
The estimate \hyperref[eq:5.48]{(48)} gives at once
\begin{equation*}
\log | P (z) | = \sum_ {n} \log \left| E _ {h} \left(\frac {z}{a _ {n}}\right) \right| \leq (2 h + 1) | z | ^ {h + 1} \sum_ {n} | a _ {n} | ^ {- h - 1}
\end{equation*}
and it follows that $P(z)$ is at most of order $h + 1$.

For the opposite inequality assume that $f(z)$ is of finite order $\lambda$ and let $h$ be the largest integer $\leq \lambda$. Then $h + 1 > \lambda$, and we have to prove, first of all, that $\sum_{n} |a_n|^{-h-1}$ converges. It is for this proof that Jensen's formula is needed.

Let us denote by $\nu(\rho)$ the number of zeros $a_{n}$ with $|a_{n}|\leq\rho$. In order to find an upper bound for $\nu(\rho)$ we apply \hyperref[eq:5.44]{(44)} with $2\rho$ in the place of $\rho$ and omit the terms $\log(2\rho/|a_{n}|)$ with $|a_{n}|\geq\rho$. We find
\begin{equation*}
\nu (\rho) \log 2 \leq \frac {1}{2 \pi} \int_ {0} ^ {2 \pi} \log | f (2 \rho e ^ {i \theta}) | d \theta - \log | f (0) |.\tag{52}\label{eq:5.52}
\end{equation*}
In view of \hyperref[eq:5.47]{(47)} it follows that $\lim_{\rho \to \infty} \nu(\rho) \rho^{-\lambda - \varepsilon} = 0$ for every $\varepsilon > 0$.

We assume now that the zeros $a_{n}$ are ordered according to absolute values: $|a_{1}| \leq |a_{2}| \leq \cdots \leq |a_{n}| \leq \cdots$. Then it is clear that $n \leq \nu(|a_{n}|)$, and from a certain n on we must have, for instance,
\begin{equation*}
n \leq \nu (| a _ {n} |) <   | a _ {n} | ^ {\lambda + \varepsilon}.
\end{equation*}
According to this inequality the series $\sum_{n}|a_n|^{-h - 1}$ has the majorant
\begin{equation*}
\sum_ {n} n ^ {- \frac {h + 1}{\lambda + \varepsilon}},
\end{equation*}
and if we choose $\varepsilon$ so that $\lambda + \varepsilon < h + 1$ the majorant converges. We have thus proved that $f(z)$ can be written in the form \hyperref[eq:5.46]{(46)} where so far $g(z)$ is only known to be entire.

It remains to prove that $g(z)$ is a polynomial of degree $\leq h$. For this purpose it is easiest to use the Poisson-Jensen formula. If the operation $(\partial/\partial x) - i(\partial/\partial y)$ is applied to both sides of the identity \hyperref[eq:5.45]{(45)}, we obtain
\begin{equation*}
\begin{array}{r l} \frac {f ^ {\prime} (z)}{f (z)} = & \sum_ {1} ^ {\nu (\rho)} (z - a _ {n}) ^ {- 1} + \sum_ {1} ^ {\nu (\rho)} \bar {a} _ {n} (\rho^ {2} - \bar {a} _ {n} z) ^ {- 1} \\ & + \frac {1}{2 \pi} \int_ {0} ^ {2 \pi} 2 \rho e ^ {i \theta} (\rho e ^ {i \theta} - z) ^ {- 2} \log | f (\rho e ^ {i \theta}) | d \theta . \end{array}
\end{equation*}
On differentiating h times with respect to z this yields
\begin{equation*}
D ^ {(h)} \frac {f ^ {\prime} (z)}{f (z)} = - h! \sum_ {1} ^ {\nu (\rho)} (a _ {n} - z) ^ {- h - 1} + h! \sum_ {1} ^ {\nu (\rho)} \bar {a} _ {n} ^ {h + 1} (\rho^ {2} - \bar {a} _ {n} z) ^ {- h - 1}\tag{53}\label{eq:5.53}
\end{equation*}
% Original PDF physical page 227
\sourcepage{227}
\begin{equation*}
+ (h + 1)! \frac {1}{2 \pi} \int_ {0} ^ {2 \pi} 2 \rho e ^ {i \theta} (\rho e ^ {i \theta} - z) ^ {- h - 2} \log | f (\rho e ^ {i \theta}) | d \theta .
\end{equation*}
It is our intention to let $\rho$ tend to $\infty$. In order to estimate the integral in \hyperref[eq:5.53]{(53)} we observe first that
\begin{equation*}
\int_ {0} ^ {2 \pi} \rho e ^ {i \theta} (\rho e ^ {i \theta} - z) ^ {- h - 2} d \theta = 0.
\end{equation*}
Therefore nothing changes if we subtract $\log M(\rho)$ from $\log |f|$. If $\rho > 2|z|$ it follows that the last term in \hyperref[eq:5.53]{(53)} has a modulus at most equal to
\begin{equation*}
(h + 1)! 2 ^ {h + 3} \rho^ {- h - 1} \frac {1}{2 \pi} \int_ {0} ^ {2 \pi} \log \frac {M (\rho)}{| f (\rho e ^ {i \theta}) |} d \theta ,
\end{equation*}
for $\log \frac{M(\rho)}{|f(\rho e^{i\theta})|} \geq 0$. But
\begin{equation*}
\frac {1}{2 \pi} \int_ {0} ^ {2 \pi} \log | f | d \theta \geq \log | f (0) |
\end{equation*}
by Jensen's formula, and $\rho^{-h-1} \log M(\rho) \to 0$ since $\lambda < h + 1$. We conclude that the integral in \hyperref[eq:5.53]{(53)} tends to 0.

As for the second sum in \hyperref[eq:5.53]{(53)}, the same preliminary inequality $\rho > 2|z|$ together with $|a_n| \leq \rho$ makes each term absolutely less than $(2 / \rho)^{h + 1}$, and the whole sum has modulus at most $2^{h + 1} \nu(\rho) \rho^{-h - 1}$. We have already proved that this tends to 0. Therefore we obtain
\begin{equation*}
D ^ {(h)} \frac {f ^ {\prime} (z)}{f (z)} = - h! \sum_ {n = 1} ^ {\infty} (a _ {n} - z) ^ {- h - 1}.\tag{54}\label{eq:5.54}
\end{equation*}
Writing $f(z) = e^{g(z)}P(z)$ we find
\begin{equation*}
g ^ {(h + 1)} (z) = D ^ {(h)} \frac {f ^ {\prime}}{f} - D ^ {(h)} \frac {P ^ {\prime}}{P}.
\end{equation*}
However, by Weierstrass's theorem the quantity $D^{(h)}(P'/P)$ can be found by separate differentiation of each factor, and in this way we obtain precisely the right-hand member of \hyperref[eq:5.54]{(54)}. Consequently, $g^{(h+1)}(z)$ is identically zero, and $g(z)$ must be a polynomial of degree $\leq h$. We have proved Theorem~\hyperref[thm:5.8]{8}.

The theorem is a factorization theorem for entire functions\index{Entire function}\index{Function!entire} of finite order $\lambda$. If $\lambda$ is not an integer, the genus h, and thereby the form of the product, is uniquely determined. If the order is integral, there is an ambiguity.

The following impressive corollary shows the strength of Hadamard's theorem\index{Hadamard's theorem}:

\begin{originalstatement}{Corollary}{}\label{cor:5.unnumbered-p227-b15}
An entire function of fractional order assumes every finite value infinitely many times.
\end{originalstatement}

% Original PDF physical page 228
\sourcepage{228}
It is clear that f and f - a have the same order for any constant a. Therefore we need only show that f has infinitely many zeros. If f has only a finite number of zeros we can divide by a polynomial and obtain a function of the same order without zeros. By the theorem it must be of the form $e^{g}$ where g is a polynomial. But it is evident that the order of $e^{g}$ is exactly the degree of g, and hence an integer. The contradiction proves the corollary.

\subsection*{EXERCISES}
\originalnumber{1}{exercise:5.3.2.9.1}
1. The characterization of the genus given in the first paragraph of Sec.~\hyperref[sec:5.3.2]{3.2} is not literally the same as the definition in Sec.~\hyperref[sec:5.2.3]{2.3}. Supply the reasoning necessary to see that the conditions are equivalent.

\originalnumber{2}{exercise:5.3.2.9.2}
2. Assume that $f(z)$ has genus zero so that
\begin{equation*}
f (z) = z ^ {m} \prod_ {n} \left(1 - \frac {z}{a _ {n}}\right).
\end{equation*}
Compare $f(z)$ with
\begin{equation*}
g (z) = z ^ {m} \prod_ {n} \left(1 - \frac {z}{| a _ {n} |}\right)
\end{equation*}
and show that the maximum modulus $\max_{|z|=r}|f(z)|$ is $\leq$ the maximum modulus of $g$, and that the minimum modulus of $f$ is $\geq$ the minimum modulus of $g$.

\setcounter{section}{3}\section{The Riemann Zeta Function}\index{Function!zeta|(}\index{Riemann zeta function|(}\index{Function!zeta}\index{Riemann zeta function}\label{sec:5.4}
The series $\sum_{n=1}^{\infty} n^{-\sigma}$ converges uniformly for all real $\sigma$ greater than or equal to a fixed $\sigma_0 > 1$. It is a majorant of the series
\begin{equation*}
\zeta (s) = \sum_ {n = 1} ^ {\infty} n ^ {- s} \quad (s = \sigma + i t),\tag{55}\label{eq:5.55}
\end{equation*}
which therefore represents an analytic function of s in the half plane Re s > 1 (see Sec.~\hyperref[sec:5.1.1]{1.1}, Ex. \hyperref[exercise:5.1.1.1.2]{2}; the notation $s = \sigma + it$ is traditional in this context).

The function $\zeta(s)$ is known as Riemann's $\zeta$ -function. It plays a central role in the applications of complex analysis to number theory. It would lead us too far astray to develop even a few of these applications in this book, but we can and will acquaint the reader with some of the more elementary properties of the $\zeta$ -function.


% Original PDF physical page 229
\index{Entire function|)}\index{Function!entire|)}\index{Hadamard's theorem|)}\sourcepage{229}
\setcounter{section}{4}\setcounter{subsection}{0}\subsection{The Product Development}\label{sec:5.4.1}
The number-theoretic properties of $\zeta(s)$ are inherent in the following connection between the $\zeta$-function and the ascending sequence of primes $p_1, p_2, \ldots, p_n, \ldots$.

\begin{originalstatement}{Theorem}{9}\label{thm:5.9}
For $\sigma = \operatorname{Re}s > 1$,
\begin{equation*}
\frac {1}{\zeta (s)} = \prod_ {n = 1} ^ {\infty} (1 - p _ {n} ^ {- s}).\tag{56}\label{eq:5.56}
\end{equation*}
\end{originalstatement}
According to Theorem~\hyperref[thm:5.6]{6} the infinite product converges uniformly for $\sigma \geq \sigma_0 > 1$ if the same is true of the series $\sum_{1}^{\infty}|p_n^{-s}| = \sum_{1}^{\infty}p_n^{-\sigma}$. Since the latter is obtained by omitting terms of $\sum_{1}^{\infty}n^{-\sigma}$, its uniform convergence for $\sigma \geq \sigma_0$ is obvious.

Under the assumption $\sigma > 1$ it is seen at once that
\begin{equation*}
\zeta (s) (1 - 2 ^ {- s}) = \sum n ^ {- s} - \sum (2 n) ^ {- s} = \sum m ^ {- s}
\end{equation*}
where m runs through the odd integers. By the same reasoning
\begin{equation*}
\zeta (s) (1 - 2 ^ {- s}) (1 - 3 ^ {- s}) = \sum m ^ {- s}
\end{equation*}
where this time m runs through all integers that are neither divisible by 2 nor by 3. More generally,
\begin{equation*}
\zeta (s) \left(1 - 2 ^ {- s}\right) \left(1 - 3 ^ {- s}\right) \cdot \cdot \cdot \left(1 - p _ {N} ^ {- s}\right) = \sum m ^ {- s},\tag{57}\label{eq:5.57}
\end{equation*}
the sum of the right being over all integers that contain none of the prime factors 2, 3, . . . , $p_N$. The first term in the sum is 1, and the next is $p_{N+1}^{-s}$. Therefore, the sum of all the terms except the first tends to zero as $N \to \infty$, and we conclude that
\begin{equation*}
\lim _ {N \rightarrow \infty} \zeta (s) \prod_ {n = 1} ^ {N} (1 - p _ {n} ^ {- s}) = 1.
\end{equation*}
This proves the theorem.

We have taken for granted that there are infinitely many primes. Actually, the reasoning can be used to prove this fact. For if $p_{N}$ were the largest prime, \hyperref[eq:5.57]{(57)} would become
\begin{equation*}
\zeta (s) (1 - 2 ^ {- s}) (1 - 3 ^ {- s}) \cdot \cdot \cdot (1 - p _ {N} ^ {- s}) = 1
\end{equation*}
and it would follow that $\zeta(\sigma)$ has a finite limit when $\sigma \to 1$. This contra%
% Original PDF physical page 230
\sourcepage{230}%
dicts the divergence of $\sum_{1}^{\infty} n^{-1}$.

\setcounter{section}{4}\setcounter{subsection}{1}\subsection{\texorpdfstring{Extension of $\zeta(s)$ to the Whole Plane}{Extension of ζ(s) to the Whole Plane}}\label{sec:5.4.2}
Recall that
\begin{equation*}
\Gamma (s) = \int_ {0} ^ {\infty} x ^ {s - 1} e ^ {- x} d x
\end{equation*}
for $\sigma > 1$ (Sec.~\hyperref[sec:5.2.5]{2.5}, \hyperref[eq:5.42]{(42)}). On replacing $x$ by $nx$ in the integral, we obtain
\begin{equation*}
n ^ {- s} \Gamma (s) = \int_ {0} ^ {\infty} x ^ {s - 1} e ^ {- n x} d x,
\end{equation*}
and summation with respect to $n$ leads to
\begin{equation*}
\zeta (s) \Gamma (s) = \int_ {0} ^ {\infty} \frac {x ^ {s - 1}}{e ^ {x} - 1} d x.\tag{58}\label{eq:5.58}
\end{equation*}
Because $\sigma > 1$ the integral is absolutely convergent at both ends, and this justifies the interchange of integration and summation. We recall that $x^{s-1}$ is unambiguously defined as $e^{(s-1)\log x}$.

Figure~\hyperref[fig:5-1]{5-1} shows two infinite paths, C and $C_{n}$, both beginning and ending near the positive real axis. For the moment we are interested only in C; its precise shape is irrelevant, as long as the radius r of the circle about the origin is $<2\pi$.

\begin{originalstatement}{Theorem}{10}\label{thm:5.10}
For $\sigma > 1$,
\begin{equation*}
\zeta (s) = - \frac {\Gamma (1 - s)}{2 \pi i} \int_ {C} \frac {(- z) ^ {s - 1}}{e ^ {z} - 1} d z\tag{59}\label{eq:5.59}
\end{equation*}
where $(-z)^{s-1}$ is defined on the complement of the positive real axis as $e^{(s-1)\log(-z)}$ with $-\pi < \operatorname{Im} \log(-z) < \pi$.

\end{originalstatement}
\bookfigure{230}{1}{FIG. 5-1.}

% Original PDF physical page 231
\sourcepage{231}
The integral is obviously convergent. By Cauchy's theorem its value does not depend on the shape of C as long as C does not enclose any multiples of $2\pi i$. In particular, we are free to let r tend to zero. It is readily seen that the integral over the circle tends to zero with r. In the limit we are left with an integral back and forth along the positive real axis. On the upper edge $(-z)^{s-1} = x^{s-1}e^{-(s-1)\pi i}$ and on the lower edge $(-z)^{s-1} = x^{s-1}e^{(s-1)\pi i}$. We obtain
\begin{equation*}
\begin{array}{r l} \int_ {C} \frac {(- z) ^ {s - 1}}{e ^ {z} - 1} d z & = - \int_ {0} ^ {\infty} \frac {x ^ {s - 1} e ^ {- (s - 1) \pi i}}{e ^ {x} - 1} d x + \int_ {0} ^ {\infty} \frac {x ^ {s - 1} e ^ {(s - 1) \pi i}}{e ^ {x} - 1} d x \\ & = 2 i \sin (s - 1) \pi \zeta (s) \Gamma (s). \end{array}
\end{equation*}
Because $\sin (s - 1)\pi = -\sin s\pi$ and $\Gamma (s)\Gamma (1 - s) = \pi /\sin s\pi$ (Sec.~\hyperref[sec:5.2.4]{2.4}, \hyperref[eq:5.30]{(30)}), this implies \hyperref[eq:5.59]{(59)}.

The importance of the formula \hyperref[eq:5.59]{(59)} lies in the fact that the right-hand side is defined and meromorphic for all values of s, so the formula can be used to extend $\zeta(s)$ to a meromorphic function in the whole plane. It is indeed quite obvious that the integral in \hyperref[eq:5.59]{(59)} is an entire function of s, while $\Gamma(1 - s)$ is meromorphic with poles at $s = 1, 2, \ldots$. Because $\zeta(s)$ is already known to be analytic for $\sigma > 1$, the poles at the integers $n \geq 2$ must cancel against zeros of the integral. At s = 1, $-\Gamma(1 - s)$ has a simple pole with the residue 1, as seen for instance by Sec.~\hyperref[sec:5.2.4]{2.4}, \hyperref[eq:5.29]{(29)}. On the other hand,
\begin{equation*}
\frac {1}{2 \pi i} \int_ {C} \frac {d z}{e ^ {z} - 1} = 1
\end{equation*}
by residues, so $\zeta(s)$ has the residue 1. We formulate the result as a corollary.

\begin{originalstatement}{Corollary}{}\label{cor:5.unnumbered-p231-b6}
The $\zeta$-function can be extended to a meromorphic function in the whole plane whose only pole is a simple pole at $s = 1$ with the residue 1.
\end{originalstatement}
The values $\zeta(-n)$ at the negative integers and zero can be evaluated explicitly. Recall the expansion (Sec.~\hyperref[sec:5.1.3]{1.3}, Ex. \hyperref[exercise:5.1.3.3.4]{4})
\begin{equation*}
\frac {1}{e ^ {z} - 1} = \frac {1}{z} - \frac {1}{2} + \sum_ {1} ^ {\infty} (- 1) ^ {k - 1} \frac {B _ {k}}{(2 k) !} z ^ {2 k - 1}.\tag{60}\label{eq:5.60}
\end{equation*}
From \hyperref[eq:5.59]{(59)}
\begin{equation*}
\zeta (- n) = (- 1) ^ {n} \frac {n !}{2 \pi i} \int_ {C} \frac {z ^ {- n - 1}}{e ^ {z} - 1} d z.
\end{equation*}
Hence $\zeta(-n)$ is equal to $(-1)^{n} n!$ times the coefficient of $z^{n}$ in \hyperref[eq:5.60]{(60)}, and
% Original PDF physical page 232
\sourcepage{232}
we can read off the following values: $\zeta(0) = -1/2$, $\zeta(-2m) = 0$, and $\zeta(-2m + 1) = (-1)^{m}B_{m}/2m$ for positive integers m. The points -2m are called the trivial zeros of the $\zeta$ -function.

\setcounter{section}{4}\setcounter{subsection}{2}\subsection{The Functional Equation}\label{sec:5.4.3}
In the half plane $\sigma > 1$ the $\zeta$ -function is given explicitly by the series \hyperref[eq:5.55]{(55)}, and it is therefore subject to the estimate $|\zeta(s)| \leq \zeta(\sigma)$. Riemann recognized that there is a rather simple relationship between $\zeta(s)$ and $\zeta(1 - s)$. As a consequence, one has good control of the behavior of the $\zeta$ -function also in the half plane $\sigma < 0$.

We shall reproduce one of the standard proofs of the functional equation\index{Functional equation|(}\index{Functional equation}, as it is commonly called.

\begin{originalstatement}{Theorem}{11}\label{thm:5.11}
\begin{equation*}
\zeta (s) = 2 ^ {s} \pi^ {s - 1} \sin \frac {\pi s}{2} \Gamma (1 - s) \zeta (1 - s).\tag{61}\label{eq:5.61}
\end{equation*}
\end{originalstatement}
For the proof we make use of the path $C_{n}$ in Fig.~\hyperref[fig:5-1]{5-1}; we assume that the square part lies on the lines $t = \pm(2n + 1)\pi$ and $\sigma = \pm(2n + 1)\pi$. The cycle $C_{n} - C$ has winding number one about the points $\pm2m\pi i$ with $m = 1, \ldots, n$. At these points the function $(-z)^{s-1}/(e^{z} - 1)$ has simple poles with residues $(\mp2m\pi i)^{s-1}$. It follows that
\begin{equation*}
\begin{array}{r l} \frac {1}{2 \pi i} \int_ {C _ {n} - C} \frac {(- z) ^ {s - 1}}{e ^ {z} - 1} d z & = \sum_ {m = 1} ^ {n} [ (- 2 m \pi i) ^ {s - 1} + (2 m \pi i) ^ {s - 1} ] \\ & = 2 \sum_ {m = 1} ^ {n} (2 m \pi) ^ {s - 1} \sin \frac {\pi s}{2}. \end{array}\tag{62}\label{eq:5.62}
\end{equation*}
We divide $C_n$ into $C_n' + C_n''$ where $C_n'$ is the part on the square and $C_n''$ the part outside the square. It is easy to see that $|e^z - 1|$ is bounded below on $C_n'$ by a fixed positive constant, independent of $n$, while $|(-z)^{s-1}|$ is bounded by a multiple of $n^{\sigma-1}$. The length of $C_n'$ is of the order of $n$, and we find that
\begin{equation*}
\left| \int_ {C _ {n} ^ {\prime}} \frac {(- z) ^ {s - 1}}{e ^ {z} - 1} d z \right| \leq A n ^ {\sigma}
\end{equation*}
for some constant $A$. If $\sigma < 0$, the integral over $C_n'$ will thus tend to zero as $n \to \infty$, and the same is of course true of the integral over $C_n''$. Therefore, the integral over $C_n - C$ will tend to the integral over $-C$, and by Theorem~\hyperref[thm:5.10]{10} the left-hand side of \hyperref[eq:5.62]{(62)} tends to $\zeta(s)/\Gamma(1 - s)$.

Under the same condition on $\sigma$ the series $\sum_{1} m^{s-1}$ converges to
% Original PDF physical page 233
\sourcepage{233}
$\zeta(1 - s)$, and the limit of the right-hand side of \hyperref[eq:5.62]{(62)} is a multiple of $\zeta(1 - s)$. The equality of the limits leads directly to the equation \hyperref[eq:5.61]{(61)}, which is thereby proved for all s with $\sigma < 0$. But two meromorphic functions which agree on a nonempty open set are identical. Hence \hyperref[eq:5.61]{(61)} is true for all s.

There are equivalent forms of the functional equation\index{Functional equation}. For instance, if we use the identity $\Gamma(s)\Gamma(1 - s) = \pi/\sin\pi s$ \hyperref[eq:5.61]{(61)} implies
\begin{equation*}
\zeta (1 - s) = 2 ^ {1 - s} \pi^ {- s} \cos \frac {\pi s}{2} \Gamma (s) \zeta (s).\tag{63}\label{eq:5.63}
\end{equation*}
The content of Theorem~\hyperref[thm:5.11]{11} can also be expressed in the following form:

\begin{originalstatement}{Corollary}{}\label{cor:5.unnumbered-p233-b4}
The function
\begin{equation*}
\xi (s) = \frac {1}{2} s (1 - s) \pi^ {- s / 2} \Gamma (s / 2) \zeta (s)
\end{equation*}
is entire and satisfies $\xi(s) = \xi(1 - s)$.

\end{originalstatement}
It is evident that $\xi(s)$ is entire, for the factor 1 - s offsets the pole of $\zeta(s)$, and the poles of $\Gamma(s/2)$ cancel against the trivial zeros of $\zeta(s)$. By use of \hyperref[eq:5.63]{(63)} the assertion $\xi(s) = \xi(1 - s)$ translates to
\begin{equation*}
\begin{array}{r l} \pi^ {- s / 2} \Gamma (s / 2) \zeta (s) & = \pi^ {(s - 1) / 2} \Gamma \left(\frac {1 - s}{2}\right) \zeta (1 - s) \\ & = 2 ^ {1 - s} \pi^ {- (s + 1) / 2} \Gamma (s) \Gamma \left(\frac {1 - s}{2}\right) \cos \frac {\pi s}{2}, \end{array}
\end{equation*}
which is the same as
\begin{equation*}
\cos \frac {\pi s}{2} \Gamma (s) \Gamma \left(\frac {1 - s}{2}\right) = 2 ^ {s - 1} \pi^ {1 / 2} \Gamma \left(\frac {s}{2}\right).
\end{equation*}
Because of the relation
\begin{equation*}
\Gamma \left(\frac {1 - s}{2}\right) \Gamma \left(\frac {1 + s}{2}\right) = \pi / \cos \frac {\pi s}{2}
\end{equation*}
the last equation is equivalent to
\begin{equation*}
\pi^ {1 / 2} \Gamma (s) = 2 ^ {s - 1} \Gamma \left(\frac {s}{2}\right) \Gamma \left(\frac {1 + s}{2}\right),
\end{equation*}
and this is nothing else than Legendre's duplication formula (Sec.~\hyperref[sec:5.2.4]{2.4}, \hyperref[eq:5.32]{(32)}). The corollary is proved.

What is the order of $\xi(s)?$ Because $\xi(s) = \xi(1 - s)$ it is sufficient to estimate $|\xi(s)|$ for $\sigma \geq \frac{1}{2}$. It is an easy consequence of Stirling's formula (Sec.~\hyperref[sec:5.2.5]{2.5}, \hyperref[eq:5.37]{(37)}) that $\log |\Gamma(s/2)| \leq A |s| \log |s|$ for some constant $A$ and
% Original PDF physical page 234
\index{Functional equation|)}\sourcepage{234}
large $|s|$, and this estimate is precise for real values of s. Therefore, if we can show that $|\zeta(s)|$ is relatively small when $\sigma \geq \frac{1}{2}$, it will follow that the order is equal to 1.

We use the standard notation $[x]$ for the largest integer $\leq x$. Assume first that $\sigma > 1$. The reader will have no difficulty verifying the following computation:
\begin{equation*}
\begin{aligned}
\int_N^\infty[x]x^{-s-1}\,dx
&=\sum_N^\infty n\int_n^{n+1}x^{-s-1}\,dx\\
&=s^{-1}\sum_N^\infty n(n^{-s}-(n+1)^{-s})\\
&=s^{-1}\left[N^{-s+1}+\sum_{N+1}^\infty n^{-s}\right].
\end{aligned}
\end{equation*}
It follows that
\begin{equation*}
\zeta (s) = \sum_ {1} ^ {N} n ^ {- s} + \frac {1}{s - 1} N ^ {1 - s} - s \int_ {N} ^ {\infty} (x - [ x ]) x ^ {- s - 1} d x.\tag{64}\label{eq:5.64}
\end{equation*}
So far this is proved for $\sigma > 1$, but the integral on the right converges for $\sigma > 0$, and the equality will therefore remain valid for $\sigma > 0$; incidentally, \hyperref[eq:5.64]{(64)} exhibits the pole at s = 1 with residue 1.

If $\sigma \geq \frac{1}{2}$ \hyperref[eq:5.64]{(64)} yields an estimate of the form
\begin{equation*}
| \zeta (s) | \leq N + A | N | ^ {- 1 / 2} | s |
\end{equation*}
valid for large $|s|$ with A independent of s and N. By choosing N as the integer closest to $|s|^{2/3}$, we find that $|\zeta(s)|$ is bounded by a constant times $|s|^{2/3}$. Therefore this factor does not influence the order.

\setcounter{section}{4}\setcounter{subsection}{3}\subsection{The Zeros of the Zeta Function}\index{Function!zeta}\label{sec:5.4.4}
It follows from the product development \hyperref[eq:5.56]{(56)} that $\zeta(s)$ has no zeros in the half plane $\sigma > 1$. With this information the functional equation implies that the only zeros in the half plane $\sigma < 0$ are the trivial ones. In other words, all nontrivial zeros lie in the so-called critical strip $0 \leq \sigma \leq 1$. The famous Riemann conjecture, which has neither been proved nor disproved, asserts that all nontrivial zeros lie on the critical line $\sigma = \frac{1}{2}$. It is not too difficult to prove that there are no zeros on $\sigma = 1$ and $\sigma = 0$. It is known that asymptotically more than one third of the zeros lie on the critical line. \hyperlink{fn:234:dagger}{\textsuperscript{\ensuremath{\dagger}}}

Let $N(T)$ be the number of zeros with $0 \leq t \leq T$. For the information of the reader we state without proof that
\begin{equation*}
N (T) = \frac {T}{2 \pi} \log \frac {T}{2 \pi} - \frac {T}{2 \pi} + O (\log T).
\end{equation*}
\begingroup\renewcommand{\thefootnote}{}\footnotetext{\hypertarget{fn:234:dagger}{}\textsuperscript{\ensuremath{\dagger}} Proved by Norman Levinson in 1975.}\endgroup

% Original PDF physical page 235
\index{Function!zeta|)}\index{Riemann zeta function|)}\sourcepage{235}
\setcounter{section}{4}\section{Normal Families}\index{Normal family|(}\index{Normal family}\label{sec:5.5}
In Chap.~\hyperref[chap:3]{3}, Sec.~\hyperref[sec:3.1]{1} we have already familiarized the reader with the idea of regarding a function as a point in a space. In principle there is thus no difference between a set of points and a set of functions. In order to make a clear distinction we shall nevertheless prefer to speak of families of functions, and usually we assume that all functions in a family are defined on the same set.

We are primarily interested in families of analytic functions, defined in a fixed region. Important examples are the families of bounded analytic functions, of functions which do not take the same value twice, etc. The aim is to study convergence properties within such families.

\setcounter{section}{5}\setcounter{subsection}{0}\subsection{Equicontinuity}\index{Equicontinuity|(}\index{Equicontinuity}\label{sec:5.5.1}
Although analytic functions are our main concern, it is expedient to choose a more general starting point. It turns out that our basic theorems are valid, and equally easy to prove, for families of functions with values in any metric space.

As a basic assumption we shall let $\mathfrak{F}$ denote a family of functions $f$, defined in a fixed region $\Omega$ of the complex plane, and with values in a metric space S. As in Chap.~\hyperref[chap:3]{3}, Sec.~\hyperref[sec:3.1]{1}, the distance function in S will be denoted by d.

We are interested in the convergence of sequences $\{f_{n}\}$ formed by functions in F. There is no particular reason to expect a sequence $\{f_{n}\}$ to be convergent; on the contrary, it is perhaps more likely that we run into the opposite extreme of a sequence that does not possess a single convergent subsequence. In many situations the latter possibility is a serious disadvantage, and the purpose of the considerations that follow is to find conditions which rule out this kind of behavior.

Let us review the definition of continuity of a function f with values in a metric space. By definition, f is continuous at $z_{0}$ if to every $\varepsilon > 0$ there exists a $\delta > 0$ such that $d(f(z), f(z_{0})) < \varepsilon$ as soon as $|z - z_{0}| < \delta$. We recall that f is said to be uniformly continuous if we can choose $\delta$ independent of $z_{0}$. But in the case of a family of functions there is another relevant kind of uniformity, namely, whether we can choose $\delta$ independent of f. We choose to require both, and are thus led to the following definition:

\begin{originalstatement}{Definition}{1}\label{def:5.1}
The functions in a family $\mathfrak{F}$ are said to be equicontinuous on a set $E \subset \Omega$ if and only if, for each $\varepsilon > 0$, there exists a $\delta > 0$ such that $d(f(z), f(z_0)) < \varepsilon$ whenever $|z - z_0| < \delta$ and $z_0, z \in E$, simultaneously for all functions $f \in \mathfrak{F}$.
\end{originalstatement}
Observe that, with this definition, each f in an equicontinuous family


% Original PDF physical page 236
\sourcepage{236}
is itself uniformly continuous on $E$.

We return now to the question of convergent subsequences. Our second definition serves to characterize families with a regular behavior:

\begin{originalstatement}{Definition}{2}\label{def:5.2}
A family $\mathfrak{F}$ is said to be normal in $\Omega$ if every sequence $\{f_{n}\}$ of functions $f_{n} \in \mathfrak{F}$ contains a subsequence which converges uniformly on every compact subset of $\Omega$.
\end{originalstatement}
This definition does not require the limit functions of the convergent subsequences to be members of F.

\setcounter{section}{5}\setcounter{subsection}{1}\subsection{Normality and Compactness}\label{sec:5.5.2}
The reader cannot fail to have noticed the close similarity between normality and the Bolzano-Weierstrass property (Chap.~\hyperref[chap:3]{3}, Theorem~\hyperref[thm:3.7]{7}). To make it more than a similarity we need to define a distance on the space of functions on $\Omega$ with values in S, and convergence with respect to this distance function should mean precisely the same as uniform convergence on compact sets.

For this purpose we need, first of all, an exhaustion of $\Omega$ by an increasing sequence of compact sets $E_{k} \subset \Omega$. By this we mean that every compact subset E of $\Omega$ shall be contained in an $E_{k}$. The construction is possible in many ways: To be specific, let $E_{k}$ consist of all points in $\Omega$ at distance $\leq k$ from the origin, and at distance $\geq 1/k$ from the boundary $\partial\Omega$. It is clear that each $E_{k}$ is bounded and closed, hence compact. Any compact set $E \subset \Omega$ is bounded and at positive distance from $\partial\Omega$; therefore it is contained in an $E_{k}$.

Let f and g be any two functions on $\Omega$ with values in S. We shall define a distance $\rho(f,g)$ between these functions, not to be confused with the distances $d(f(z),g(z))$ between their values. To do so we first replace d by the distance function
\begin{equation*}
\delta (a, b) = \frac {d (a , b)}{1 + d (a , b)}
\end{equation*}
which also satisfies the triangle inequality and has the advantage of being bounded (Chap.~\hyperref[chap:3]{3}, Sec.~\hyperref[sec:3.1.2]{1.2}, Ex. \hyperref[exercise:3.1.2.1.1]{1}). Next, we set
\begin{equation*}
\delta_ {k} (f, g) = \sup _ {z \in E _ {k}} \delta (f (z), g (z))
\end{equation*}
which may be described as the distance between f and g on $E_{k}$. Finally, we agree on the definition
\begin{equation*}
\rho (f, g) = \sum_ {k = 1} ^ {\infty} \delta_ {k} (f, g) 2 ^ {- k}.\tag{65}\label{eq:5.65}
\end{equation*}
It is trivial to verify that $\rho(f,g)$ is finite and satisfies all the conditions for a distance function (Chap.~\hyperref[chap:3]{3}, Sec.~\hyperref[sec:3.1.2]{1.2}).


% Original PDF physical page 237
\index{Equicontinuity|)}\sourcepage{237}
The distance $\rho(f,g)$ has the property we were looking for. Suppose first that $f_{n}\to f$ in the sense of the $\rho$ -distance. For large n we have then $\rho(f_{n},f)<\varepsilon$ and consequently, by \hyperref[eq:5.65]{(65)}, $\delta_{k}(f_{n},f)<2^{k}\varepsilon$. But this implies that $f_{n}\to f$ uniformly on $E_{k}$, first with respect to the $\delta$ -metric, but hence also with respect to the d-metric. Since every compact E is contained in an $E_{k}$ it follows that the convergence is uniform on E.

Conversely, suppose that $f_{n}$ converges uniformly to $f$ on every compact set. Then $\delta_k(f_n,f) \to 0$ for every $k$, and because the series $\Sigma \delta_k(f_n,f)2^{-k}$ has a convergent majorant with terms independent of $n$ it follows readily (as in Weierstrass's $M$ test) that $\rho(f_n,f) \to 0$.

We have shown that convergence with respect to the distance $\rho$ is equivalent to convergence on compact sets. So far we did not assume S to be complete, but if it is, it follows easily that the space of all functions with values in S is complete as a metric space with the distance $\rho$.

It can be said with some justification that the metric we have introduced is arbitrary and artificial. However, from what we have proved it follows that the open sets are independent of the choices involved in the construction. In other words, the topology has an intrinsic meaning, tailored to the needs of the theory of analytic functions.

We now recall the Bolzano-Weierstrass theorem, according to which a metric space is compact if and only if every infinite sequence has a convergent subsequence (Chap.~\hyperref[chap:3]{3}, Theorem~\hyperref[thm:3.7]{7}). The theorem is applied to the set F, equipped with the distance $\rho$, and we conclude that F is compact if and only if F is normal, and if the limit functions are themselves in F. On the other hand, if F is normal, so is its closure $F^{-}$. Therefore we obtain the following characterization of normal families\index{Normal family}:

\begin{originalstatement}{Theorem}{12}\label{thm:5.12}
A family $\mathfrak{F}$ is normal if and only if its closure $\mathfrak{F}^{-}$ with respect to the distance function \hyperref[eq:5.65]{(65)} is compact.
\end{originalstatement}
It is also customary to say that $\mathfrak{F}$ is relatively compact if $\mathfrak{F}^{-}$ is compact. Thus, normal and relatively compact families are the same.

We shall now relate the notion of normal families to total boundedness. If $\mathfrak{F}$ is normal, then $\mathfrak{F}^{-}$ is compact, and according to Chap.~\hyperref[chap:3]{3}, Theorem~\hyperref[thm:3.6]{6}, $\mathfrak{F}^{-}$ is totally bounded, and so is consequently $\mathfrak{F}$ (see the footnote on p.~\hyperlink{orig-p61}{61}). By definition, this means that to every $\epsilon > 0$ there exist a finite number of functions $f_{1}, \ldots, f_{n} \in \mathfrak{F}$ such that every $f \in \mathfrak{F}$ satisfies $\rho(f, f_{j}) > \varepsilon$ for some $f_{j}$. Conversely, if $\mathfrak{F}$ is totally bounded, so is $\mathfrak{F}^{-}$. If $S$ is known to be complete, then $\mathfrak{F}^{-}$ is also complete, and hence compact. In other words, if $S$ is complete, then $\mathfrak{F}$ is normal if and only if it is totally bounded.

The following theorem serves to state the condition of total boundedness in terms of the original metric on S rather than in terms of the auxiliary metric $\rho$.


% Original PDF physical page 238
\sourcepage{238}
\begin{originalstatement}{Theorem}{13}\label{thm:5.13}
The family $\mathfrak{F}$ is totally bounded if and only if to every compact set $E\subset\Omega$ and every $\varepsilon>0$ it is possible to find $f_{1},\ldots,f_{n}\in\mathfrak{F}$ such that every $f\in\mathfrak{F}$ satisfies $d(f,f_{j})<\varepsilon$ on $E$ for some $f_{j}$.
\end{originalstatement}
If $\mathfrak{F}$ is totally bounded there exist $f_{1},\ldots ,f_{n}$ such that, for any $f\in \mathfrak{F}$, $\rho (f,f_j) < \varepsilon$ for some $f_{j}$. By \hyperref[eq:5.65]{(65)} this implies $\delta_k(f,f_j) < 2^k\varepsilon$, or $\delta (f,f_j) < 2^k\varepsilon$ on $E_{k}$. If we fix $k$ beforehand, we can thus make $\delta (f,f_j)$ arbitrarily small on $E_{k}$, and the same is then true of $d(f,f_j)$. This proves that the condition is necessary.

To prove the sufficiency we choose $k_0$ so that $2^{-k_0} < \varepsilon / 2$. By assumption we can find $f_1, \ldots, f_n$ such that any $f \in \mathfrak{F}$ satisfies one of the inequalities $\delta(f, f_j) \leq d(f, f_j) < \varepsilon / 2k_0$ on $E_{k_0}$. It follows that $\delta_k(f, f_j) < \varepsilon / 2k_0$ for $k \leq k_0$, while trivially $\delta_k(f, f_j) < 1$ for $k > k_0$. From \hyperref[eq:5.65]{(65)} we obtain
\begin{equation*}
\rho (f, f _ {j}) <   k _ {0} \left(\varepsilon / 2 k _ {0}\right) + 2 ^ {- k _ {0} - 1} + 2 ^ {- k _ {0} - 2} + \dots = \varepsilon / 2 + 2 ^ {- k _ {0}} <   \varepsilon ,
\end{equation*}
which is precisely what we wanted to prove.

\setcounter{section}{5}\setcounter{subsection}{2}\subsection{Arzela's Theorem}\label{sec:5.5.3}
We shall now study the relationship between Definition~\hyperref[def:5.1]{1} and Definition~\hyperref[def:5.2]{2}. The connection is established by a famous and extremely useful theorem known as Arzela's theorem (or the Arzela-Ascoli theorem\index{Arzela-Ascoli theorem|(}\index{Arzela-Ascoli theorem}).

\begin{originalstatement}{Theorem}{14}\label{thm:5.14}
A family F of continuous functions with values in a metric space S is normal in the region $\Omega$ of the complex plane if and only if
(i) $\mathfrak{F}$ is equicontinuous on every compact set $E\subset \Omega$

(ii) for any $z \in \Omega$ the values $f(z), f \in \mathfrak{F}$, lie in a compact subset of $S$.

\end{originalstatement}
We give two proofs of the necessity of (i). Assume that $\mathfrak{F}$ is normal and determine $f_{1},\ldots ,f_{n}$ as in Theorem~\hyperref[thm:5.13]{13}. Because each of these functions is uniformly continuous on $E$ we can find a $\delta >0$ such that $d(f_j(z),f_j(z_0)) < \varepsilon$ for $z,z_0\in E$, $|z - z_0| < \delta$, $j = 1,\dots ,n$. For any given $f\in \mathfrak{F}$ and corresponding $f_{j}$ we obtain
\begin{equation*}
\begin{aligned}
d(f(z),f(z_0))\leq{}&d(f(z),f_j(z))+d(f_j(z),f_j(z_0))\\
&+d(f_j(z_0),f(z_0))<3\varepsilon
\end{aligned}
\end{equation*}
and (i) is proved.

Less elegantly, but without use of Theorem~\hyperref[thm:5.13]{13}, a proof can be given as follows: If $\mathfrak{F}$ fails to be equicontinuous on $E$ there exists an $\varepsilon > 0$, sequences of points $z_n, z_n' \in E$, and functions $f_n \in \mathfrak{F}$ such that $|z_n - z_n'| \to 0$ while $d(f_n(z_n), f_n(z_n')) \geq \varepsilon$ for all $n$. Because $E$ is compact we can choose subsequences of $\{z_n\}$ and $\{z_n'\}$ which converge to a common limit $z'' \in E$, and because $\mathfrak{F}$ is normal there exists a subsequence of $\{f_n\}$ which converges uniformly on $E$. It is clear that we may choose all three sub%
% Original PDF physical page 239
\sourcepage{239}%
sequences to have the same subscripts $n_k$. The limit function $f$ of $\{f_{n_k}\}$ is uniformly continuous on $E$. Hence we can find $k$ such that the distances from $f_{n_k}(z_{n_k})$ to $f(z_{n_k})$, from $f(z_{n_k})$ to $f(z_{n_k}')$, and from $f(z_{n_k}')$ to $f_{n_k}(z_{n_k}')$ are all $< \varepsilon / 3$. It follows that $d(f_{n_k}(z_{n_k}), f_{n_k}(z_{n_k}')) < \varepsilon$, contrary to the assumption that $d(f_n(z_n), f_n(z_n')) \geq \varepsilon$ for all $n$.

To prove the necessity of (ii) we show that the closure of the set formed by the values $f(z), f \in \mathfrak{F}$, is compact. Let $\{w_n\}$ be a sequence in this closure. To each $w_n$ we determine $f_n \in \mathfrak{F}$ so that $d(f_n(z), w_n) < 1/n$. By normality there exists a convergent subsequence $\{f_{n_k}(z)\}$, and the sequence $\{w_{n_k}\}$ converges to the same value.

The sufficiency of (i) together with (ii) is proved by Cantor\index{Cantor G@Cantor, G.}'s famous diagonal process. We observe first that there exists an everywhere dense sequence of points $\zeta_{k}$ in $\Omega$, for instance the points with rational coordinates. From the sequence $\{f_n\}$ we are going to extract a subsequence which converges at all points $\zeta_{k}$. To find a subsequence which converges at one given point is always possible because of condition (ii). We can therefore find an array of subscripts
\begin{equation*}
\begin{array}{l} n _ {1 1} <   n _ {1 2} <   \cdot \cdot \cdot <   n _ {1 j} <   \cdot \cdot \cdot \\ n _ {2 1} <   n _ {2 2} <   \cdot \cdot \cdot <   n _ {2 j} <   \cdot \cdot \cdot \\ \cdot \cdot \cdot \\ n _ {k 1} <   n _ {k 2} <   \cdot \cdot \cdot <   n _ {k j} <   \cdot \cdot \cdot \\ \cdot \cdot \cdot \\ \end{array}\tag{66}\label{eq:5.66}
\end{equation*}
such that each row is contained in the preceding one, and such that $\lim_{j\to\infty}f_{n_{kj}}(\zeta_{k})$ exists for all k. The diagonal sequence $\{n_{jj}\}$ is strictly increasing, and it is ultimately a subsequence of each row in \hyperref[eq:5.66]{(66)}. Hence $\{f_{n_{jj}}\}$ is a subsequence of $\{f_{n}\}$ which converges at all points $\zeta_{k}$. For convenience we replace the notation $n_{jj}$ by $n_{j}$.

Consider now a compact set $E \subset \Omega$ and assume that $\mathfrak{F}$ is equicontinuous on $E$. We shall show that $\{f_{n_j}\}$ converges uniformly on $E$. Given $\varepsilon > 0$ we choose $\delta > 0$ such that, for $z, z' \in E$ and $f \in \mathfrak{F}, |z - z'| < \delta$ implies $d(f(z), f(z')) < \varepsilon/3$. Because $E$ is compact, it can be covered by a finite number of $\delta/2$-neighborhoods. We select a point $\zeta_k$ from each of these neighborhoods. There exists an $i_0$ such that $i,j > i_0$ implies $d(f_{n_i}(\zeta_k), f_{n_j}(\zeta_k)) < \varepsilon/3$ for all these $\zeta_k$. For each $z \in E$ one of the $\zeta_k$ is within distance $\delta$ from $z$; hence $d(f_{n_i}(z), f_{n_i}(\zeta_k)) < \varepsilon/3$, $d(f_{n_j}(z), f_{n_j}(\zeta_k)) < \varepsilon/3$. The three inequalities yield $d(f_{n_i}(z), f_{n_j}(z)) < \varepsilon$. Because all values $f(z)$ belong to a compact and consequently complete subset of $S$ it follows that $\{f_{n_j}\}$ is uniformly convergent on $E$.

\setcounter{section}{5}\setcounter{subsection}{3}\subsection{Families of Analytic Functions}\label{sec:5.5.4}
Analytic functions have their values in $\mathbf{C}$, the finite complex plane. In order to apply the preceding
% Original PDF physical page 240
\index{Arzela-Ascoli theorem|)}\sourcepage{240}
considerations to families of analytic functions it is therefore natural to choose $S = \mathbf{C}$ with the usual euclidean distance.

The compact subsets of $\mathbf{C}$ are the bounded and closed sets. For this reason condition (ii) in Arzela's theorem is fulfilled if and only if the values $f(z)$ are bounded for each $z \in \Omega$, with a bound that may depend on $z$. Suppose now that condition (i) is also satisfied. For a given $z_0 \in \Omega$ determine $\rho$ so that the closed disk $|z - z_0| \leq \rho$ is contained in $\Omega$. Then $\mathfrak{F}$, the given family of functions, is equicontinuous on the closed disk. If in the definition of equicontinuity $\delta(<\rho)$ corresponds to $\varepsilon$, and if $|f(z_0)| \leq M$ for all $f \in \mathfrak{F}$, then $|f(z)| \leq M + \varepsilon$ in $|z - z_0| < \delta$. Because any compact set can be covered by a finite number of neighborhoods with this property, it follows that the functions are uniformly bounded on every compact set, the bound depending on the set. According to Arzela's theorem this is true for all normal families\index{Normal family} of complex-valued functions.

For analytic functions this condition is also sufficient.

\begin{originalstatement}{Theorem}{15}\label{thm:5.15}
A family $\mathfrak{F}$ of analytic functions is normal with respect to $\mathbf{C}$ if and only if the functions in $\mathfrak{F}$ are uniformly bounded on every compact set.
\end{originalstatement}
To prove the sufficiency we prove equicontinuity. Let C be the boundary of a closed disk in $\Omega$, of radius r. If $z, z_{0}$ are inside C we obtain by Cauchy's integral theorem
\begin{equation*}
\begin{array}{r l} f (z) - f (z _ {0}) & = \frac {1}{2 \pi i} \int_ {C} \left(\frac {1}{\zeta - z} - \frac {1}{\zeta - z _ {0}}\right) f (\zeta) d \zeta \\ & = \frac {z - z _ {0}}{2 \pi i} \int_ {C} \frac {f (\zeta) d \zeta}{(\zeta - z) (\zeta - z _ {0})}. \end{array}
\end{equation*}
If $|f| \leq M$ on $C$, and if we restrict $z$ and $z_0$ to the smaller concentric disk of radius $r/2$, it follows that
\begin{equation*}
\left| f (z) - f \left(z _ {0}\right) \right| \leq \frac {4 M \left| z - z _ {0} \right|}{r}.\tag{67}\label{eq:5.67}
\end{equation*}
This proves equicontinuity on the smaller disk.

Let $E$ be a compact set in $\Omega$. Each point of $E$ is the center of a disk with radius $r$, as above. The open disks of radius $r/4$ form an open covering of $E$. We select a finite subcovering and denote the corresponding centers, radii, and bounds by $\zeta_k$, $r_k$, $M_k$; let $r$ be the smallest of the $r_k$ and $M$ the largest of the $M_k$. For a given $\varepsilon > 0$ let $\delta$ be the smaller of $r/4$ and $\varepsilon r/4M$. If $|z - z_0| < \delta$ and $|z_0 - \zeta_k| < r_k/4$ it follows that $|z - \zeta_k| < \delta + r_k/4 \leq r_k/2$. Hence \hyperref[eq:5.67]{(67)} is applicable and we find $|f(z) - f(z_0)| \leq 4M_k\delta/r_k \leq 4M\delta/r \leq \varepsilon$ as desired.

In view of Theorem~\hyperref[thm:5.15]{15} we may abandon the term “normal with
% Original PDF physical page 241
\sourcepage{241}
respect to C" which has no historic justification. If a family has the property of the theorem, we say instead that it is locally bounded. Indeed, if the family is bounded in a neighborhood of each point, then it is obviously bounded on every compact set. The theorem tells us that every sequence has a subsequence which converges uniformly on compact sets if and only if it is locally bounded.

An interesting feature is that local boundedness is inherited by the derivatives.

\begin{originalstatement}{Theorem}{16}\label{thm:5.16}
A locally bounded family\index{Locally bounded family} of analytic functions has locally bounded derivatives.
\end{originalstatement}
This follows at once by the Cauchy representation of the derivative. If C is the boundary of a closed disk in $\Omega$, of radius r, then
\begin{equation*}
f ^ {\prime} (z) = \frac {1}{2 \pi i} \int_ {C} \frac {f (\zeta) d \zeta}{(\zeta - z) ^ {2}}.
\end{equation*}
Hence $|f'(z)| \leq 4M/r$ in the concentric disk of radius r/2 (M is the bound of $|f|$ on C). We see that the $f'$ are indeed locally bounded.

What is true of the first derivatives is of course also true of higher derivatives.

\setcounter{section}{5}\setcounter{subsection}{4}\subsection{The Classical Definition}\label{sec:5.5.5}
If a sequence tends to $\infty$ there is no great scattering of values, and it may well be argued that for the purposes of normal families\index{Normal family} such a sequence should be regarded as convergent. This is the classical point of view, and we shall restyle our definition to conform with traditional usage.

\begin{originalstatement}{Definition}{3}\label{def:5.3}
A family of analytic functions in a region $\Omega$ is said to be normal if every sequence contains either a subsequence that converges uniformly on every compact set $E \subset \Omega$, or a subsequence that tends uniformly to $\infty$ on every compact set.
\end{originalstatement}
We shall show that this definition agrees with Definition~\hyperref[def:5.2]{2} if we take S to be the Riemann sphere. If that is what we do, then we can also allow $\infty$ as a possible value, which means that we may consider families of meromorphic functions. There is no need to rephrase the definition so that it covers normal families of meromorphic functions, for Definition~\hyperref[def:5.2]{2} applies without change.

It is necessary, however, to prove a lemma which extends Weierstrass's and Hurwitz\index{Hurwitz A@Hurwitz, A.|(}\index{Hurwitz A@Hurwitz, A.}'s theorems to meromorphic functions (Theorems~\hyperref[thm:5.1]{1} and \hyperref[thm:5.2]{2}).


% Original PDF physical page 242
\sourcepage{242}
\begin{originalstatement}{Lemma}{}\label{lem:5.unnumbered-p242-b0}
If a sequence of meromorphic functions converges in the sense of spherical distance, uniformly on every compact set, then the limit function is meromorphic or identically equal to $\infty$.
If a sequence of analytic functions converges in the same sense, then the limit function is either analytic or identically equal to $\infty$.
\end{originalstatement}

Suppose $f(z) = \lim_{n\to \infty}f_n(z)$ in the sense of the lemma. We know that $f(z)$ is continuous in the spherical metric. If $f(z_0)\neq \infty$, then $f(z)$ is bounded in a neighborhood of $z_0$, and for large $n$ the functions $f_{n}$ are $\neq \infty$ in the same neighborhood. It follows by the ordinary form of Weierstrass's theorem that $f(z)$ is analytic in a neighborhood of $z_0$. If $f(z_0) = \infty$ we consider the reciprocal $1 / f(z)$ which is the limit of $1 / f_{n}(z)$ in the spherical sense. We conclude that $1 / f(z)$ is analytic near $z_0$, and hence $f(z)$ is meromorphic. If the $f_{n}$ are analytic and the second case occurs, then $1 / f$ must be identically zero by virtue of Hurwitz\index{Hurwitz A@Hurwitz, A.}'s theorem, and $f$ is identically $\infty$.

The lemma makes it clear that Definition~\hyperref[def:5.3]{3} is nothing other than Definition~\hyperref[def:5.2]{2} applied to the spherical metric.

It is not true that the derivatives of a normal family\index{Normal family} form a normal family. For instance, consider the family formed by the functions $f_{n}=n(z^{2}-n)$ in the whole plane. This family is normal, for it is clear that $f_{n}\to\infty$ uniformly on every compact set. Nevertheless, the derivatives $f_{n}^{\prime}=2nz$ do not form a normal family, for $f_{n}^{\prime}(z)$ tends to $\infty$ for $z\neq0$ and to 0 for z=0.

By Arzela's theorem a family of meromorphic functions is normal if and only if it is equicontinuous on compact sets, for condition (ii) is now trivially fulfilled. The equicontinuity can be replaced by a boundedness condition. We have indeed:

\begin{originalstatement}{Theorem}{17}\label{thm:5.17}
A family of analytic or meromorphic functions $f$ is normal in the classical sense if and only if the expressions
\begin{equation*}
\rho (f) = \frac {2 | f ^ {\prime} (z) |}{1 + | f (z) | ^ {2}}\tag{58}\label{eq:5.58-p242-b7}
\end{equation*}
are locally bounded.\hyperlink{fn:242:dagger}{\textsuperscript{\ensuremath{\dagger}}}
\end{originalstatement}

The geometric meaning of the quantity $\rho(f)$ is rather evident. Indeed, by use of the formula in Chap.~\hyperref[chap:1]{1}, Sec.~\hyperref[sec:1.2.4]{2.4}
\begin{equation*}
d (f (z _ {1}), f (z _ {2})) = \frac {2 | f (z _ {1}) - f (z _ {2}) |}{[ (1 + | f (z _ {1}) | ^ {2}) (1 + | f (z _ {2}) | ^ {2}) ] ^ {\frac {1}{2}}}
\end{equation*}
\begingroup\renewcommand{\thefootnote}{}\footnotetext{\hypertarget{fn:242:dagger}{}\textsuperscript{\ensuremath{\dagger}} This theorem is due to F. Marty\index{Marty F@Marty, F.|hyperpagenote}.}\endgroup

% Original PDF physical page 243
\index{Hurwitz A@Hurwitz, A.|)}\sourcepage{243}
it is readily seen that f followed by stereographic projection maps an arc $\gamma$ on an image with length
\begin{equation*}
\int_ {\gamma} \rho (f (z)) | d z |.
\end{equation*}
If $\rho(f) \leq M$ on the line segment between $z_1$ and $z_2$ we conclude that $d(f(z_1), f(z_2)) \leq M |z_1 - z_2|$, and this immediately proves the equicontinuity when $\rho(f)$ is locally bounded.

To prove the necessity we remark first that $\rho(f) = \rho(1/f)$ as a simple calculation shows. Assume that the family $\mathfrak{F}$ of meromorphic functions is normal, but that the $\rho(f)$ fail to be bounded on a compact set $E$. Consider a sequence of $f_n \in \mathfrak{F}$ such that the maximum of $\rho(f_n)$ on $E$ tends to $\infty$. Let $f$ denote the limit function of a convergent subsequence $\{f_{n_k}\}$. Around each point of $E$ we can find a small closed disk, contained in $\Omega$, on which either $f$ or $1/f$ is analytic. If $f$ is analytic it is bounded on the closed disk, and it follows by the spherical convergence that $\{f_{n_k}\}$ has no poles in the disk as soon as $k$ is sufficiently large. We can then use Weierstrass's theorem (Theorem~\hyperref[thm:5.1]{1}) to conclude that $\rho(f_{n_k}) \to \rho(f)$, uniformly on a slightly smaller disk. Since $\rho(f)$ is continuous it follows that $\rho(f_{n_k})$ is bounded on the smaller disk. If $1/f$ is analytic the same proof applies to $\rho(1/f_{n_k})$, which is the same as $\rho(f_{n_k})$. In conclusion, since $E$ is compact it can be covered by a finite number of the smaller disks, and we find that the $\rho(f_{n_k})$ are bounded on $E$, contrary to assumption. The contradiction completes the proof of the theorem.

\subsection*{EXERCISES}
\originalnumber{1}{exercise:5.5.5.10.1}
1. Prove that in any region $\Omega$ the family of analytic functions with positive real part is normal. Under what added condition is it locally bounded? Hint: Consider the functions $e^{-f}$.

\originalnumber{2}{exercise:5.5.5.10.2}
2. Show that the functions $z^{n}$, $n$ a nonnegative integer, form a normal family\index{Normal family} in $|z| < 1$, also in $|z| > 1$, but not in any region that contains a point on the unit circle.

\originalnumber{3}{exercise:5.5.5.10.3}
3. If $f(z)$ is analytic in the whole plane, show that the family formed by all functions $f(kz)$ with constant $k$ is normal in the annulus $r_{1} < |z| < r_{2}$ if and only if $f$ is a polynomial.

\originalnumber{4}{exercise:5.5.5.10.4}
4. If the family $\mathfrak{F}$ of analytic (or meromorphic) functions is not normal in $\Omega$, show that there exists a point $z_0$ such that $\mathfrak{F}$ is not normal in any neighborhood of $z_0$. Hint: A compactness argument.


% Original PDF physical page 244
\index{Normal family|)}\sourcepage{244}
